%=====================================================================
% MATERIAL 03 - INSTALACOES ELETRICAS PREDIAIS E INDUSTRIAIS
% Versao modular: evita duplicacao e divergencia entre os materiais
% Atualizacao tecnica: agosto de 2026
%=====================================================================
\documentclass[11pt,aspectratio=169]{beamer}

\usepackage[utf8]{inputenc}
\usepackage[T1]{fontenc}
\usepackage[brazil]{babel}
\usepackage{lmodern}
\usepackage{amsmath,amssymb,bm}
\usepackage{siunitx}
\usepackage{booktabs,array,tabularx,multirow}
\usepackage[most]{tcolorbox}
\usepackage{tikz}
\usepackage{pgfplots}
\usepackage[european,american currents]{circuitikz}
\usepackage{ragged2e}
\usepackage{microtype}
\usepackage{xcolor}
\usetikzlibrary{arrows.meta,positioning,calc,fit,shapes.geometric,patterns,decorations.pathmorphing,decorations.markings}
\pgfplotsset{compat=1.18}
\sisetup{locale=DE,per-mode=symbol,detect-all}

\newcommand{\Disciplina}{Instalações Elétricas Prediais e Industriais}
\newcommand{\TituloCurto}{Material 03}
\newcommand{\sen}{\operatorname{sen}}
\newcommand{\tg}{\operatorname{tg}}
\newcommand{\fase}[1]{\textcolor{corPrim}{\textbf{#1}}}
\newcommand{\alerta}[1]{\textcolor{corAt}{\textbf{#1}}}

%=====================================================================
% ESTILO DA CASA -- TEMPLATE BEAMER (padrao da colecao)
% Arquivo central de identidade visual para todos os materiais.
% Azul-petroleo + ambar | tcolorbox | TikZ/pgfplots/circuitikz
%=====================================================================

% Dependencias minimas do proprio estilo.
\RequirePackage{xcolor}
\RequirePackage{tikz}
\usetikzlibrary{calc}
\RequirePackage{tcolorbox}
\tcbuselibrary{skins,breakable}
\RequirePackage{listings}
\RequirePackage{xparse}
\RequirePackage{etoolbox}

\makeatletter

%====================== PALETA OFICIAL ================================
\definecolor{corPrim}{RGB}{13,48,82}      % azul-petroleo
\definecolor{corAccent}{RGB}{222,138,28}  % ambar
\definecolor{corDef}{RGB}{31,111,168}     % azul medio
\definecolor{corNorma}{RGB}{0,121,107}    % verde-petroleo
\definecolor{corEx}{RGB}{46,125,50}       % verde
\definecolor{corFix}{RGB}{204,102,0}      % laranja
\definecolor{corAt}{RGB}{176,58,46}       % vermelho
\definecolor{codBg}{RGB}{252,250,245}
\definecolor{codKw}{RGB}{13,48,82}
\definecolor{codStr}{RGB}{46,125,50}
\definecolor{codCom}{RGB}{120,130,138}

% Compatibilidade com nomes usados historicamente nos materiais.
\colorlet{Petroleo}{corPrim}
\colorlet{PetroleoEscuro}{corPrim!82!black}
\colorlet{PetroleoMedio}{corDef!82!corPrim}
\colorlet{PetroleoClaro}{corDef!7}
\colorlet{Ambar}{corAccent}
\colorlet{AmbarEscuro}{corAccent!72!black}
\colorlet{AmbarClaro}{corAccent!16}
\colorlet{AzulPetroleo}{corPrim}
\colorlet{AzulPetroleoEscuro}{corPrim!82!black}
\colorlet{AzulPetroleoClaro}{corDef!7}
\colorlet{AzulMedio}{corDef}
\colorlet{AzulClaro}{corDef!7}
\colorlet{CinzaEscuro}{black!82}
\colorlet{CinzaClaro}{black!5}
\colorlet{CinzaSuave}{black!3}
\colorlet{CinzaMedio}{black!55}
\colorlet{CinzaTexto}{black!82}
\colorlet{CinzaBox}{black!3}
\colorlet{Verde}{corEx}
\colorlet{Vermelho}{corAt}
\colorlet{Turquesa}{corNorma!70!corDef}
\colorlet{Roxo}{corDef!58!corAt}
\colorlet{ifmagreen}{corPrim}
\colorlet{ifmagreen2}{corDef}
\colorlet{ifmalight}{corDef!7}
\colorlet{accent}{corAccent}
\colorlet{warn}{corAt}
\colorlet{ok}{corEx}
\colorlet{bluex}{corDef}
\colorlet{grayx}{black!58}

% Variantes em minusculas presentes em materiais antigos.
\colorlet{azulPetroleo}{corPrim}
\colorlet{azulPetroleoEscuro}{corPrim!82!black}
\colorlet{azulPetroleoClaro}{corDef!7}
\colorlet{azulMedio}{corDef}
\colorlet{azulClaro}{corDef!7}
\colorlet{azulEscuro}{corPrim!82!black}
\colorlet{azul}{corPrim}
\colorlet{azul2}{corDef}
\colorlet{ambar}{corAccent}
\colorlet{ambarClaro}{corAccent!16}
\colorlet{cinza}{black!5}
\colorlet{cinzaClaro}{black!5}
\colorlet{cinzaMedio}{black!55}
\colorlet{cinzaEscuro}{black!82}
\colorlet{cinzaTexto}{black!82}
\colorlet{verde}{corEx}
\colorlet{verdeclaro}{corEx!10}
\colorlet{vermelho}{corAt}
\colorlet{vermelhoclaro}{corAt!9}
\colorlet{laranja}{corFix}
\colorlet{roxo}{corDef!58!corAt}

%====================== TEMA BEAMER ==================================
\usetheme{default}
\usefonttheme{professionalfonts}
\setbeamercolor{normal text}{fg=black!85,bg=white}
\setbeamercolor{structure}{fg=corPrim}
\setbeamercolor{frametitle}{bg=corPrim,fg=white}
\setbeamercolor{footsec}{bg=corPrim,fg=white!92!corPrim}
\setbeamercolor{alerted text}{fg=corAccent!85!black}
\setbeamerfont{frametitle}{series=\bfseries,size=\large}
\setbeamerfont{footline}{size=\tiny}
\setbeamertemplate{navigation symbols}{}
\setbeamertemplate{caption}{\raggedright\insertcaption\par}
\setbeamersize{text margin left=0.55cm,text margin right=0.55cm}

\setbeamertemplate{itemize item}{\raise0.35ex\hbox{\color{corAccent}\rule{0.62ex}{0.62ex}}}
\setbeamertemplate{itemize subitem}{\color{corDef}\raise0.5pt\hbox{--}}
\setbeamertemplate{itemize subsubitem}{\color{corPrim}\tiny$\circ$}
\setbeamertemplate{enumerate item}{\color{corPrim}\textbf{\insertenumlabel.}}
\setbeamertemplate{section in toc}{%
  \leavevmode\textcolor{corAccent}{\textbf{\inserttocsectionnumber.}}~\inserttocsection\par\vskip1.5pt}
\setbeamertemplate{subsection in toc}{%
  \leavevmode\leftskip=0.95cm\textcolor{corDef}{\textbullet}~\inserttocsubsection\par}

%====================== TITULO DO SLIDE ==============================
\setbeamertemplate{frametitle}{%
  \nointerlineskip%
  \begin{beamercolorbox}[wd=\paperwidth,ht=2.7ex,dp=1.3ex,leftskip=0.55cm,rightskip=0.55cm]{frametitle}%
    \usebeamerfont{frametitle}\insertframetitle%
  \end{beamercolorbox}%
  {\color{corAccent}\rule{\paperwidth}{1.3pt}}\par\vskip2pt%
}

%====================== RODAPE DUPLO ================================
\newcommand{\CasaRodapeEsq}{%
  \ifcsname Disciplina\endcsname\Disciplina\else\insertshorttitle\fi}
\newcommand{\CasaRodapeDir}{%
  \ifcsname TituloCurto\endcsname\TituloCurto\else\insertshorttitle\fi}
\newcommand{\CasaMetaCapa}{%
  \ifcsname Disciplina\endcsname
    \if\relax\detokenize\expandafter{\Disciplina}\relax
    \else
      \Disciplina\ \textbullet\ \insertdate
    \fi
  \else
    \insertdate
  \fi}
\setbeamertemplate{footline}{%
  \leavevmode%
  \hbox{%
    \begin{beamercolorbox}[wd=.5\paperwidth,ht=2.4ex,dp=1ex,leftskip=0.5cm]{footsec}%
      \usebeamerfont{footline}\CasaRodapeEsq%
    \end{beamercolorbox}%
    \begin{beamercolorbox}[wd=.5\paperwidth,ht=2.4ex,dp=1ex,rightskip=0.5cm]{footsec}%
      \hfill\usebeamerfont{footline}\CasaRodapeDir\quad\textbar\quad\insertframenumber%
    \end{beamercolorbox}%
  }\vskip0pt%
}

%====================== DIVISORIAS DE SECAO =========================
\newcommand{\CasaSecNum}{\ifnum\value{section}<10 0\fi\arabic{section}}
\newcommand{\CasaDivisoriaManual}[2]{%
  \begin{frame}[plain,noframenumbering]
    \makebox[0pt][l]{%
      \raisebox{0pt}[0pt][0pt]{%
        \begin{tikzpicture}[remember picture,overlay]
          \fill[corPrim] (current page.south west) rectangle (current page.north east);
          \if\relax\detokenize{#1}\relax\else
            \node[anchor=west,text=corAccent,font=\bfseries\fontsize{40}{40}\selectfont]
              at ([xshift=1.6cm,yshift=1.1cm]current page.west) {#1};
          \fi
          \fill[corAccent] ([xshift=1.7cm,yshift=0.35cm]current page.west)
            rectangle ++(2.4,0.09);
          \node[anchor=west,text=white,font=\bfseries\LARGE,text width=0.75\paperwidth]
            at ([xshift=1.6cm,yshift=-0.5cm]current page.west) {#2};
        \end{tikzpicture}%
      }%
    }%
  \end{frame}%
}
\AtBeginSection[]{\CasaDivisoriaManual{\CasaSecNum}{\insertsection}}

% Alguns materiais antigos usam \secslide{titulo}; outros, \secslide{numero}{titulo}.
\@ifundefined{secslide}
  {\NewDocumentCommand{\secslide}{m g}{\IfNoValueTF{#2}{\CasaDivisoriaManual{}{#1}}{\CasaDivisoriaManual{#1}{#2}}}}
  {\RenewDocumentCommand{\secslide}{m g}{\IfNoValueTF{#2}{\CasaDivisoriaManual{}{#1}}{\CasaDivisoriaManual{#1}{#2}}}}
\@ifundefined{Divisoria}
  {\NewDocumentCommand{\Divisoria}{m g}{\IfNoValueTF{#2}{\CasaDivisoriaManual{}{#1}}{\CasaDivisoriaManual{#1}{#2}}}}
  {\RenewDocumentCommand{\Divisoria}{m g}{\IfNoValueTF{#2}{\CasaDivisoriaManual{}{#1}}{\CasaDivisoriaManual{#1}{#2}}}}

%====================== CAPA ========================================
\defbeamertemplate*{title page}{estilo da casa}{%
  \vspace{0.8cm}%
  {\color{white}\bfseries\huge\inserttitle\par}%
  \vspace{2mm}%
  {\color{corAccent}\rule{3.2cm}{2.2pt}\par}%
  \vspace{4mm}%
  {\color{white!85}\large\insertsubtitle\par}%
  \vspace{8mm}%
  {\color{white}\normalsize\insertauthor\par}%
  \vspace{1mm}%
  {\color{white!70}\small\CasaMetaCapa\par}%
}
\setbeamertemplate{title page}[estilo da casa]

%====================== CAIXAS =======================================
% Materiais antigos usam convencoes diferentes para o argumento opcional
% das caixas. Para preservar a semantica, caixas existentes nao sao
% redefinidas; apenas sua paleta historica e recolorida acima.
\tcbset{
  caixabase/.style={
    enhanced,sharp corners=downhill,boxrule=0.6pt,left=2mm,right=2mm,
    top=1.2mm,bottom=1.2mm,fonttitle=\bfseries\small,
    attach boxed title to top left={xshift=3mm,yshift=-2.4mm},
    boxed title style={sharp corners,boxrule=0pt},
    before skip=3pt,after skip=3pt
  }
}
\@ifundefined{caixaDef}{\newtcolorbox{caixaDef}[1][]{caixabase,colback=corDef!5,colframe=corDef,coltitle=white,colbacktitle=corDef,title=Definição,#1}}{}
\@ifundefined{caixaNorma}{\newtcolorbox{caixaNorma}[1][]{caixabase,colback=corNorma!5,colframe=corNorma,coltitle=white,colbacktitle=corNorma,title=Norma / Método,#1}}{}
\@ifundefined{caixaEx}{\newtcolorbox{caixaEx}[1][]{caixabase,colback=corEx!5,colframe=corEx,coltitle=white,colbacktitle=corEx,title=Exemplo resolvido,#1}}{}
\@ifundefined{caixaFix}{\newtcolorbox{caixaFix}[1][]{caixabase,colback=corFix!5,colframe=corFix,coltitle=white,colbacktitle=corFix,title=Fixação,#1}}{}
\@ifundefined{caixaAt}{\newtcolorbox{caixaAt}[1][]{caixabase,colback=corAt!5,colframe=corAt,coltitle=white,colbacktitle=corAt,title=Atenção,#1}}{}
\@ifundefined{caixaForm}{\newtcolorbox{caixaForm}[1][]{caixabase,colback=corPrim!4,colframe=corPrim,coltitle=white,colbacktitle=corPrim,title=Formulário,#1}}{}
\@ifundefined{caixaObs}{\newtcolorbox{caixaObs}[1][]{caixabase,colback=black!3,colframe=corPrim,coltitle=white,colbacktitle=corPrim,title=Observação,#1}}{}
\@ifundefined{caixaAlerta}{\newtcolorbox{caixaAlerta}[1][]{caixabase,colback=corAt!5,colframe=corAt,coltitle=white,colbacktitle=corAt,title=Atenção,#1}}{}
\@ifundefined{caixaProjeto}{\newtcolorbox{caixaProjeto}[1][]{caixabase,colback=corNorma!5,colframe=corNorma,coltitle=white,colbacktitle=corNorma,title=Projeto prático,#1}}{}
\@ifundefined{caixaProj}{\newtcolorbox{caixaProj}[1][]{caixabase,colback=corNorma!5,colframe=corNorma,coltitle=white,colbacktitle=corNorma,title=Projeto prático,#1}}{}

%====================== CODIGO =======================================
\lstdefinestyle{codCasa}{
  basicstyle=\ttfamily\scriptsize,
  keywordstyle=\color{codKw}\bfseries,
  stringstyle=\color{codStr},commentstyle=\color{codCom}\itshape,
  numbers=left,numberstyle=\tiny\color{codCom},numbersep=5pt,
  frame=single,rulecolor=\color{corPrim!25},backgroundcolor=\color{codBg},
  breaklines=true,showstringspaces=false,tabsize=2,columns=fullflexible
}
\lstdefinestyle{codC}{style=codCasa,language=C}
\lstdefinestyle{codPy}{style=codCasa,language=Python}
\lstdefinestyle{py}{style=codCasa,language=Python}

%====================== EXTENSOES MODULARES ==========================
% Se existir <jobname>_antes_laboratorios.tex, o arquivo e carregado uma
% unica vez imediatamente antes da secao de laboratorios. Isso permite
% ampliar um material sem duplicar o arquivo principal.
\newif\ifCasaExtensaoAntesLabsInserida
\CasaExtensaoAntesLabsInseridafalse
\AtBeginDocument{%
  \IfFileExists{\jobname_antes_laboratorios.tex}{%
    \let\CasaSectionOriginal\section
    \renewcommand{\section}[1]{%
      \ifCasaExtensaoAntesLabsInserida\else
        \ifstrequal{##1}{Laboratórios, projetos e comissionamento}{%
          \global\CasaExtensaoAntesLabsInseridatrue
          \input{\jobname_antes_laboratorios.tex}%
        }{}%
      \fi
      \CasaSectionOriginal{##1}%
    }%
  }{}%
}

\makeatother


\title{Instalações Elétricas Prediais e Industriais}
\subtitle{Projeto, dimensionamento, proteção, normas atuais e aplicações no Maranhão}
\author{Coleção de Engenharia Elétrica}
\institute{Material 03 — versão combinada modular}
\date{Atualização: agosto de 2026}

\begin{document}

{\setbeamercolor{background canvas}{bg=corPrim}
\setbeamertemplate{footline}{}
\begin{frame}[plain,noframenumbering]
  \titlepage
\end{frame}}

\begin{frame}{Estrutura desta edição}
\small
\begin{columns}[T,onlytextwidth]
\column{0.49\textwidth}
\begin{caixaNorma}[title=Predial]
\begin{itemize}
\item planta arquitetônica antes do lançamento elétrico;
\item previsão de iluminação, TUG e TUE;
\item circuitos, quadro, DR/DDR, DPS e aterramento;
\item projeto residencial completo para o Maranhão;
\item padrões atuais da Equatorial Energia.
\end{itemize}
\end{caixaNorma}
\column{0.49\textwidth}
\begin{caixaProjeto}[title=Industrial]
\begin{itemize}
\item alimentadores, QGBT, CCM e motores;
\item curto-circuito e seletividade;
\item segurança, qualidade de energia e eficiência;
\item documentação, exemplos e exercícios.
\end{itemize}
\end{caixaProjeto}
\end{columns}
\vspace{1mm}
\begin{caixaObs}
Este arquivo não mantém mais uma cópia paralela dos capítulos. Ele importa os mesmos módulos usados nos materiais Predial e Industrial, evitando conteúdo desatualizado ou contraditório.
\end{caixaObs}
\end{frame}

% Base comum
\section{Normas, escopo e ciclo de projeto}

\begin{frame}{Instalação elétrica: visão de sistema}
\small
\begin{columns}[T,onlytextwidth]
\column{0.52\textwidth}
\begin{caixaDef}
Uma instalação elétrica é um \textbf{sistema coordenado}: fonte, entrada, medição, distribuição, proteção, condutores, aterramento, cargas, comando, documentação e procedimentos de operação/manutenção.
\end{caixaDef}
\begin{itemize}
\item Predial: residência, comércio, escola, hospital, laboratório e serviços.
\item Industrial: máquinas, motores, fornos, automação, instrumentação, CCM e redes internas.
\item O objetivo não é apenas ``funcionar'': é garantir segurança, desempenho, continuidade e manutenção.
\end{itemize}
\column{0.46\textwidth}
\centering
\begin{tikzpicture}[scale=0.74,every node/.style={font=\scriptsize}]
\tikzset{b/.style={draw=corPrim,fill=corDef!7,rounded corners,minimum width=3.0cm,minimum height=0.65cm,align=center}}
\node[b] (r) at (0,3.2) {Rede / fonte local};
\node[b] (e) at (0,2.2) {Entrada e medição};
\node[b] (q) at (0,1.2) {QGBT / QD};
\node[b] (c1) at (-1.8,0.1) {Circuitos\\prediais};
\node[b] (c2) at (1.8,0.1) {Circuitos\\industriais};
\node[b] (s) at (0,-1.1) {PE + equipotencial\\+ DPS + SPDA};
\draw[-{Stealth},thick,corPrim] (r)--(e)--(q);
\draw[-{Stealth},thick,corPrim] (q)--(c1);
\draw[-{Stealth},thick,corPrim] (q)--(c2);
\draw[-{Stealth},thick,corNorma] (s)--(q);
\draw[-{Stealth},thick,corNorma] (s)--(c1);
\draw[-{Stealth},thick,corNorma] (s)--(c2);
\end{tikzpicture}
\end{columns}
\end{frame}

\begin{frame}[shrink=7]{Mapa normativo — situação em agosto de 2026}
\scriptsize
\begin{tabularx}{\textwidth}{>{\bfseries}p{2.8cm}p{3.1cm}X}
\toprule
Documento & Situação usada no material & Aplicação \\
\midrule
ABNT NBR 5410:2004 & Referência vigente adotada para BT; a própria NT.00001.EQTL Rev.09 a cita expressamente. & Instalações elétricas de baixa tensão: choque, condutores, circuitos, DR, aterramento, verificação etc. \\
ABNT NBR 14039:2021 & Referência para instalações de média tensão. & Instalações de 1,0 kV a 36,2 kV. \\
ABNT NBR 5419-1:2026 & Nova edição publicada em 2026. & Princípios e determinação da necessidade de proteção contra descargas atmosféricas; conferir as partes aplicáveis da série no projeto. \\
ABNT NBR 17019:2022 & Referência para alimentação de veículos elétricos. & Instalações fixas para recarga de VE; também referenciada pela NT.00042.EQTL Rev.04. \\
NR-10 + Portaria MTE 737/2026 & A nova redação foi publicada, mas a Portaria entra em vigor um ano após sua publicação. & Segurança em instalações e serviços com eletricidade; distinguir requisito vigente de requisito já publicado com vigência futura. \\
Série ABNT NBR IEC 61439 & Usar edição aplicável ao conjunto. & Quadros e conjuntos de manobra e comando de baixa tensão. \\
\bottomrule
\end{tabularx}
\end{frame}

\begin{frame}[shrink=8]{Equatorial Energia — referências atuais para o Maranhão}
\scriptsize
\begin{tabularx}{\textwidth}{>{\bfseries}p{2.65cm}p{2.15cm}X}
\toprule
Documento & Revisão/data & Quando usar \\
\midrule
NT.00001.EQTL & Rev.09 — 29/05/2025 & Fornecimento de energia em baixa tensão para unidade consumidora individual. É a referência principal do projeto residencial deste material. \\
NT.00004.EQTL & Rev.08 — 31/12/2025 & Empreendimentos/edificações de múltiplas unidades consumidoras — condomínios, edifícios e centros de medição. \\
NT.00020.EQTL & Rev.06 — 31/12/2025 & Conexão de micro e minigeração distribuída; formulários do Grupo B atualizados em 2026. \\
NT.00030.EQTL & Rev.03 — 31/12/2025 & Padrões construtivos das caixas de medição e proteção. \\
NT.00042.EQTL & Rev.04 — 31/12/2025 & Critérios de conexão de estações de recarga de veículos elétricos. \\
NT.00045.EQTL & Rev.01 — 29/03/2023 & Postes para padrão de entrada. \\
\bottomrule
\end{tabularx}
\begin{caixaAt}
As normas da concessionária tratam da interface com a rede e dos padrões de conexão. Elas \textbf{não substituem} o dimensionamento da instalação interna conforme ABNT e demais requisitos legais.
\end{caixaAt}
\end{frame}

\begin{frame}{Equatorial Maranhão: limites que afetam o projeto residencial}
\small
\begin{columns}[T,onlytextwidth]
\column{0.50\textwidth}
\begin{caixaNorma}[title=NT.00001.EQTL Rev.09]
A norma atende unidades individuais em BT com carga instalada de até 75 kW. Para o Maranhão:
\begin{itemize}
\item ligação monofásica urbana: 220 V, dentro dos limites previstos;
\item ligação trifásica urbana: \textbf{380/220 V, 3F+N};
\item acima dos limites e nos casos especiais, é necessária análise específica.
\end{itemize}
\end{caixaNorma}
\column{0.48\textwidth}
\begin{caixaProjeto}[title=Tabela 1 — 220/380 V]
Exemplos de enquadramento por carga instalada:
\begin{itemize}
\item até 5 kW: mono 25 A;
\item 5,1--8 kW: mono 40 A;
\item 8,1--12 kW: mono 63 A;
\item 12,1--24 kW: tri 40 A;
\item 24,1--38 kW: tri 63 A.
\end{itemize}
\end{caixaProjeto}
\end{columns}
\begin{caixaObs}
O projeto residencial completo deste material fecha em aproximadamente 19,88 kW instalados e, por isso, usa o enquadramento trifásico de 40 A como \textbf{consequência} da carga calculada.
\end{caixaObs}
\end{frame}

\begin{frame}{Projeto interno e projeto para a concessionária não são a mesma coisa}
\small
\begin{columns}[T,onlytextwidth]
\column{0.49\textwidth}
\begin{caixaNorma}[title=Unidade individual em BT]
A NT.00001.EQTL informa que, em regra, não é necessário apresentar à concessionária um projeto elétrico separado para novas instalações/reformas de unidade individual atendida em tensão secundária.
\end{caixaNorma}
\column{0.49\textwidth}
\begin{caixaAt}[title=Mas continuam obrigatórios]
\begin{itemize}
\item projeto interno tecnicamente correto;
\item padrão de entrada conforme concessionária;
\item ABNT, segurança e Corpo de Bombeiros quando aplicável;
\item normas específicas para GD, VE e EMUC;
\item responsabilidade técnica quando exigida.
\end{itemize}
\end{caixaAt}
\end{columns}
\end{frame}

\begin{frame}{NR-10: nova redação publicada, vigência futura}
\small
\begin{columns}[T,onlytextwidth]
\column{0.52\textwidth}
\begin{caixaNorma}[title=Portaria MTE nº 737/2026]
A Portaria é de 29/05/2026, foi publicada no DOU de 01/06/2026 e estabelece que sua nova redação da NR-10 \textbf{entra em vigor um ano após a publicação}.
\end{caixaNorma}
\begin{itemize}
\item integra explicitamente a gestão com o GRO/NR-1;
\item reforça choque e arco elétrico;
\item detalha segurança em projetos e atualização documental;
\item prevê estudo de energia incidente quando aplicável.
\end{itemize}
\column{0.46\textwidth}
\begin{caixaAt}[title=Como aplicar em agosto de 2026]
Não tratar a redação nova como se já estivesse vigente. No material:
\begin{itemize}
\item requisito atual é apresentado como atual;
\item texto publicado para 2027 é identificado como futuro;
\item boas práticas podem ser antecipadas sem falsa alegação normativa.
\end{itemize}
\end{caixaAt}
\end{columns}
\end{frame}

\begin{frame}{Ciclo completo de um projeto elétrico}
\small
\centering
\begin{tikzpicture}[node distance=0.32cm,every node/.style={font=\scriptsize}]
\tikzset{et/.style={draw=corPrim,fill=corDef!7,rounded corners,minimum width=2.55cm,minimum height=0.72cm,align=center}}
\node[et] (a) {1. Levantamento\\arquitetura e cargas};
\node[et,right=of a] (b) {2. Normas, risco\\e premissas};
\node[et,right=of b] (c) {3. Circuitos e\\unifilar};
\node[et,below=of c] (d) {4. Cálculos\\e proteção};
\node[et,left=of d] (e) {5. Quadros e\\infraestrutura};
\node[et,left=of e] (f) {6. Memorial, BIM\\e especificação};
\node[et,below=of f] (g) {7. Execução e\\fiscalização};
\node[et,right=of g] (h) {8. Ensaios e\\comissionamento};
\node[et,right=of h] (i) {9. As built e\\manutenção};
\foreach \u/\v in {a/b,b/c,c/d,d/e,e/f,f/g,g/h,h/i}{\draw[-{Stealth},thick,corPrim] (\u)--(\v);}
\draw[-{Stealth},thick,corAccent] (i.east) to[out=0,in=-70] node[right,align=center]{retroalimentar\\o projeto} (a.south east);
\end{tikzpicture}
\end{frame}

\begin{frame}{Entregáveis mínimos de um bom projeto}
\small
\begin{columns}[T,onlytextwidth]
\column{0.49\textwidth}
\begin{itemize}
\item Planta com pontos, rotas e identificação.
\item Quadro de cargas, demanda e balanceamento.
\item Diagramas unifilares e, quando necessário, multifilares.
\item Memorial de cálculo: corrente, cabos, queda, curto, proteção e aterramento.
\item Detalhes de quadros, BEP, DPS, SPDA e interfaces.
\end{itemize}
\column{0.49\textwidth}
\begin{itemize}
\item Especificação técnica e lista de materiais.
\item Estudos adicionais quando aplicáveis.
\item Plano de inspeção/comissionamento.
\item As built e registro de alterações.
\item ART/RRT e demais responsabilidades conforme competência profissional.
\end{itemize}
\end{columns}
\end{frame}

%=====================================================================

\section{Fundamentos, potência, energia e demanda}

\begin{frame}{Grandezas em corrente alternada}
\small
\begin{columns}[T,onlytextwidth]
\column{0.54\textwidth}
\begin{align*}
v(t)&=V_m\sen(\omega t+\theta_v)\\
i(t)&=I_m\sen(\omega t+\theta_i)\\
V&=\frac{V_m}{\sqrt{2}},\qquad I=\frac{I_m}{\sqrt{2}}\\
\varphi&=\theta_v-\theta_i
\end{align*}
\begin{itemize}
\item RMS relaciona a grandeza CA ao efeito térmico equivalente em CC.
\item Em carga indutiva, corrente tende a atrasar a tensão.
\item Em carga capacitiva, corrente tende a adiantar a tensão.
\end{itemize}
\column{0.44\textwidth}
\begin{tikzpicture}
\begin{axis}[width=5.5cm,height=3.7cm,axis lines=middle,xlabel={$\omega t$},xtick={0,90,180,270,360},xticklabels={$0$,$90$,$180$,$270$,$360$},ytick=\empty,domain=0:360,samples=120,ymin=-1.3,ymax=1.3,grid=both]
\addplot[corPrim,thick] {sin(x)};
\addplot[corAccent,thick,dashed] {0.85*sin(x-35)};
\end{axis}
\node[corPrim,font=\scriptsize] at (1.25,3.25) {$v(t)$};
\node[corAccent,font=\scriptsize] at (4.25,2.45) {$i(t)$};
\end{tikzpicture}
\end{columns}
\end{frame}

\begin{frame}{Potência ativa, reativa e aparente}
\small
\begin{columns}[T,onlytextwidth]
\column{0.50\textwidth}
\begin{align*}
S &= VI \quad [\si{VA}]\\
P &= VI\cos\varphi \quad [\si{W}]\\
Q &= VI\sen\varphi \quad [\si{var}]\\
S^2 &= P^2+Q^2,\qquad fp=\frac{P}{S}
\end{align*}
\begin{itemize}
\item $P$: potência convertida em trabalho, calor, luz etc.
\item $Q$: intercâmbio de energia reativa.
\item $S$: capacidade elétrica exigida de cabos, quadros e transformadores.
\end{itemize}
\column{0.48\textwidth}
\centering
\begin{tikzpicture}[scale=0.9]
\draw[->,thick] (0,0)--(4.1,0) node[right]{$P$};
\draw[->,thick] (4.1,0)--(4.1,2.4) node[above]{$Q$};
\draw[->,thick,corPrim] (0,0)--(4.1,2.4) node[midway,above left]{$S$};
\draw (0.8,0) arc (0:30:0.8);
\node at (1.15,0.28) {$\varphi$};
\end{tikzpicture}
\end{columns}
\end{frame}

\begin{frame}{Sistemas trifásicos}
\small
\begin{columns}[T,onlytextwidth]
\column{0.52\textwidth}
\begin{align*}
P_{3\phi}&=\sqrt3\,V_{LL}I_L fp\\
Q_{3\phi}&=\sqrt3\,V_{LL}I_L\sen\varphi\\
S_{3\phi}&=\sqrt3\,V_{LL}I_L
\end{align*}
\begin{itemize}
\item Cargas equilibradas reduzem corrente de neutro fundamental.
\item Cargas monofásicas devem ser distribuídas entre as fases.
\item Harmônicas triplas podem somar no neutro mesmo com correntes eficazes semelhantes nas fases.
\end{itemize}
\column{0.46\textwidth}
\centering
\begin{tikzpicture}[scale=0.95]
\draw[->,thick,corPrim] (0,0)--(2,0) node[right]{$V_A$};
\draw[->,thick,corDef] (0,0)--({2*cos(120)},{2*sin(120)}) node[left]{$V_B$};
\draw[->,thick,corAccent] (0,0)--({2*cos(240)},{2*sin(240)}) node[left]{$V_C$};
\draw[thin,black!25] (0,0) circle (2);
\node at (0,-2.35) {defasagem de $120^\circ$};
\end{tikzpicture}
\end{columns}
\end{frame}

\begin{frame}{Corrente de projeto}
\small
\begin{columns}[T,onlytextwidth]
\column{0.53\textwidth}
\begin{caixaDef}
A corrente de projeto $I_b$ é a corrente prevista no circuito considerando potência, tensão, rendimento, fator de potência e regime de funcionamento.
\end{caixaDef}
\begin{align*}
I_{1\phi}&=\frac{P}{V\,fp\,\eta}\\
I_{3\phi}&=\frac{P}{\sqrt3\,V_{LL}\,fp\,\eta}
\end{align*}
\column{0.45\textwidth}
\begin{caixaEx}
Motor de $\SI{15}{kW}$, $\SI{380}{V}$, $fp=0{,}85$, $\eta=0{,}92$:
\[
I_b=\frac{15000}{\sqrt3\cdot380\cdot0{,}85\cdot0{,}92}\approx\SI{29.2}{A}.
\]
Esse valor ainda não resolve partida, cabo, disjuntor ou queda de tensão.
\end{caixaEx}
\end{columns}
\end{frame}

\begin{frame}{Carga instalada, demanda e diversidade}
\small
\begin{columns}[T,onlytextwidth]
\column{0.50\textwidth}
\begin{align*}
P_{inst}&=\sum P_k\\
D_{max}&=\max_t P(t)\\
F_d&=\frac{D_{max}}{P_{inst}}\\
F_{div}&=\frac{\sum D_{max,k}}{D_{max,grupo}}
\end{align*}
\begin{itemize}
\item A entrada não deve ser dimensionada supondo simultaneidade total sem justificativa.
\item Fator de demanda deve vir de regra aplicável, histórico, processo ou critério documentado.
\end{itemize}
\column{0.48\textwidth}
\begin{tikzpicture}
\begin{axis}[width=5.7cm,height=3.6cm,xlabel={hora},ylabel={$P$ (kW)},xmin=0,xmax=24,ymin=0,ymax=100,grid=both,xtick={0,6,12,18,24}]
\addplot[corPrim,thick,smooth] coordinates {(0,18)(3,15)(6,30)(9,54)(12,71)(15,63)(18,90)(21,50)(24,26)};
\addplot[corAccent,dashed,thick] coordinates {(0,90)(24,90)};
\end{axis}
\end{tikzpicture}
\end{columns}
\end{frame}

\begin{frame}{Levantamento de cargas moderno}
\scriptsize
\begin{tabularx}{\textwidth}{>{\bfseries}p{2.6cm}p{3.1cm}X}
\toprule
Carga & Dados mínimos & O que acrescentar ao projeto \\
\midrule
Iluminação & Potência, comando, ambiente & Iluminância, emergência, automação e eficiência. \\
TUG/TUE & Potência, tensão, uso & Circuitos, DR, simultaneidade, expansão. \\
Motores & kW, regime, partida & $I_N$, inrush, torque, proteção, VFD/soft-starter. \\
TI/eletrônica & W/VA, UPS & Harmônicas, DPS, continuidade, aterramento funcional. \\
Ar-condicionado & Potência e tecnologia & Inverter, corrente máxima, diversidade e partida. \\
VE & Potência de recarga, modo & Circuito dedicado, DR, gestão de carga, expansão. \\
FV/armazenamento & Potência, inversor & Proteções CA/CC, anti-ilhamento, fluxos bidirecionais. \\
Processo industrial & Ciclo e criticidade & Continuidade, redundância, seletividade, produção. \\
\bottomrule
\end{tabularx}
\end{frame}

%=====================================================================


% Projeto predial
\section{Projeto predial de baixa tensão}

\begin{frame}{Projeto predial: sequência correta}
\small
\centering
\begin{tikzpicture}[node distance=0.32cm,every node/.style={font=\scriptsize}]
\tikzset{et/.style={draw=corPrim,fill=corDef!7,rounded corners,minimum width=2.55cm,minimum height=0.72cm,align=center}}
\node[et] (a) {1. Planta e uso\\dos ambientes};
\node[et,right=of a] (b) {2. Pontos e\\cargas};
\node[et,right=of b] (c) {3. Circuitos\\terminais};
\node[et,below=of c] (d) {4. Condutores e\\eletrodutos};
\node[et,left=of d] (e) {5. Proteções\\e aterramento};
\node[et,left=of e] (f) {6. Quadro e\\balanceamento};
\node[et,below=of f] (g) {7. Entrada da\\concessionária};
\node[et,right=of g] (h) {8. Unifilar e\\memorial};
\node[et,right=of h] (i) {9. Verificação\\e as built};
\foreach \u/\v in {a/b,b/c,c/d,d/e,e/f,f/g,g/h,h/i}{\draw[-{Stealth},thick,corPrim] (\u)--(\v);}
\end{tikzpicture}
\vspace{2mm}
\begin{caixaAt}
Não se começa o projeto escolhendo disjuntor, cabo ou padrão de entrada. Esses elementos são consequência da arquitetura, das cargas, dos critérios normativos e das condições reais de instalação.
\end{caixaAt}
\end{frame}

\begin{frame}{Da arquitetura para os pontos elétricos}
\small
\begin{columns}[T,onlytextwidth]
\column{0.50\textwidth}
\begin{caixaDef}[title=Levantamento]
Antes de lançar símbolos elétricos, identificar:
\begin{itemize}
\item dimensões e perímetros dos ambientes;
\item portas, janelas e circulação;
\item bancadas, mobiliário e equipamentos fixos;
\item áreas molhadas, externas e sujeitas a lavagem;
\item posição tecnicamente adequada do QD.
\end{itemize}
\end{caixaDef}
\column{0.48\textwidth}
\begin{caixaProjeto}[title=Depois]
\begin{itemize}
\item calcular os mínimos normativos;
\item acrescentar pontos necessários ao uso real;
\item definir TUE pelas cargas previstas;
\item posicionar interruptores junto aos acessos;
\item lançar rotas compatíveis com a construção.
\end{itemize}
\end{caixaProjeto}
\end{columns}
\end{frame}

\begin{frame}{Iluminação e tomadas: mínimo não é projeto final}
\small
\begin{columns}[T,onlytextwidth]
\column{0.49\textwidth}
\begin{caixaNorma}[title=Previsão mínima de iluminação]
Em habitações, a previsão de carga parte de:
\begin{itemize}
\item 100 VA para os primeiros 6 m$^2$ ou fração;
\item 60 VA para cada acréscimo de 4 m$^2$ inteiros ou em fração.
\end{itemize}
A quantidade real de luminárias decorre do projeto luminotécnico, não dessa potência de previsão.
\end{caixaNorma}
\column{0.49\textwidth}
\begin{caixaNorma}[title=TUG]
A quantidade mínima depende do tipo e do perímetro do ambiente. Cozinha, copa, lavanderia e área de serviço possuem critérios mais exigentes que salas, dormitórios e outros cômodos.
\end{caixaNorma}
\begin{caixaAt}
Não desenhar tomadas apenas para atingir uma contagem: conferir portas, móveis, bancadas, eletrodomésticos e acessibilidade de uso.
\end{caixaAt}
\end{columns}
\end{frame}

\begin{frame}{Divisão de circuitos: critérios defensáveis}
\small
\begin{columns}[T,onlytextwidth]
\column{0.49\textwidth}
\begin{itemize}
\item separar iluminação e tomadas quando tecnicamente adequado;
\item separar cargas de potência elevada ou uso específico;
\item prever circuitos exclusivos para TUE quando requerido;
\item limitar a consequência de uma falha a uma área razoável;
\item facilitar identificação, manutenção e expansão.
\end{itemize}
\column{0.49\textwidth}
\begin{itemize}
\item respeitar ambientes molhados e externos;
\item evitar compartilhamento indevido de neutros;
\item considerar DR/DDR e sua seletividade funcional;
\item distribuir circuitos monofásicos entre as fases;
\item revisar simultaneidade e perfil das cargas.
\end{itemize}
\end{columns}
\begin{caixaObs}
A divisão de circuitos deve aparecer de forma idêntica na planta, no quadro de cargas, no diagrama unifilar e na identificação física do QD.
\end{caixaObs}
\end{frame}

\begin{frame}{Quadro de cargas: não usar números soltos}
\scriptsize
\begin{tabularx}{\textwidth}{c p{2.25cm} c c c c c c X}
\toprule
Circ. & Descrição & Fase & $P$ & $I_b$ & Cabo & DJ & DR & Verificações \\
 & & & W/VA & A & mm$^2$ & A & mA & \\
\midrule
C1 & Iluminação social & C & 720 & 3,27 & 1,5 & 10 & -- & $I_z$, queda, curto \\
C2 & Iluminação íntima & B & 860 & 3,91 & 1,5 & 10 & -- & $I_z$, queda, curto \\
C3 & TUG áreas secas & A & 1100 & 5,00 & 2,5 & 16 & conf. & uso e extensão \\
C4 & TUG cozinha & B & 1900 & 8,64 & 2,5 & 16 & 30 & ambiente molhado \\
C7 & Chuveiro & A & 5500 & 25,0 & 6 & 32 & 30 & dedicado \\
\bottomrule
\end{tabularx}
\vspace{2mm}
\begin{caixaAt}
Os valores acima pertencem ao projeto residencial desenvolvido adiante. Eles não são uma tabela universal para copiar em qualquer residência.
\end{caixaAt}
\end{frame}

\begin{frame}{Planta elétrica: o que precisa existir no desenho}
\small
\begin{columns}[T,onlytextwidth]
\column{0.47\textwidth}
\begin{itemize}
\item paredes, portas, janelas e nomes dos ambientes;
\item cotas ou referência dimensional confiável;
\item luminárias, interruptores e comandos;
\item TUG, TUE, dados e demais sistemas;
\item número do circuito junto aos pontos/rotas;
\item QD em local acessível e coerente.
\end{itemize}
\column{0.51\textwidth}
\begin{itemize}
\item eletrodutos ou rotas com leitura construtiva;
\item seção/quantidade de condutores quando pertinente;
\item alturas e observações de montagem;
\item legenda de símbolos;
\item compatibilização com hidráulica, estrutura, HVAC e dados;
\item detalhes dos locais especiais.
\end{itemize}
\begin{caixaProjeto}
Uma planta elétrica não é um diagrama abstrato colocado sobre retângulos. Ela deve permitir compreender \textbf{onde} instalar, \textbf{o que} instalar e \textbf{a qual circuito} cada ponto pertence.
\end{caixaProjeto}
\end{columns}
\end{frame}

\begin{frame}{Unifilar: fechar a lógica elétrica}
\small
\begin{columns}[T,onlytextwidth]
\column{0.48\textwidth}
\begin{itemize}
\item origem e tensão de fornecimento;
\item medição e proteção de entrada;
\item alimentador e número de condutores;
\item DG e barramentos do QD;
\item DPS e conexão ao sistema de equipotencialização;
\item DR/DDR/RCBO conforme circuitos protegidos.
\end{itemize}
\column{0.50\textwidth}
\begin{itemize}
\item disjuntores terminais;
\item fases atribuídas aos circuitos;
\item neutros corretamente segregados a jusante dos DR;
\item barramento PE contínuo e identificado;
\item seções e correntes nominais;
\item identificação coerente com planta e quadro de cargas.
\end{itemize}
\end{columns}
\begin{caixaAt}
O símbolo de terra isolado não substitui o desenho do PE, do BEP e da equipotencialização. Em projeto didático, esconder essa topologia ensina errado.
\end{caixaAt}
\end{frame}

%=====================================================================


% Dimensionamento e proteção comuns
\section{Condutores, infraestrutura e queda de tensão}

\begin{frame}{Dimensionamento: não existe um único critério}
\small
\begin{columns}[T,onlytextwidth]
\column{0.52\textwidth}
\begin{caixaNorma}[title=Fluxo de decisão]
A seção final deve atender simultaneamente a:
\begin{itemize}
\item seção mínima;
\item capacidade de condução de corrente;
\item queda de tensão;
\item curto-circuito/estabilidade térmica;
\item proteção contra choque/desligamento automático;
\item requisitos mecânicos e de instalação.
\end{itemize}
\end{caixaNorma}
\column{0.46\textwidth}
\centering
\begin{tikzpicture}[node distance=0.25cm,every node/.style={font=\scriptsize}]
\tikzset{b/.style={draw=corPrim,fill=corDef!7,rounded corners,align=center,minimum width=3.6cm,minimum height=0.58cm}}
\node[b] (a) {$I_b$ e regime};
\node[b,below=of a] (b) {Método de instalação};
\node[b,below=of b] (c) {Fatores de correção};
\node[b,below=of c] (d) {$I_b\le I_n\le I_z$};
\node[b,below=of d] (e) {Queda de tensão};
\node[b,below=of e] (f) {Curto + PE + desligamento};
\foreach \u/\v in {a/b,b/c,c/d,d/e,e/f}{\draw[-{Stealth},thick,corPrim] (\u)--(\v);}
\end{tikzpicture}
\end{columns}
\end{frame}

\begin{frame}{Coordenação carga--proteção--condutor}
\small
\begin{columns}[T,onlytextwidth]
\column{0.51\textwidth}
\begin{align*}
I_b &\le I_n \le I_z\\
I_2 &\le 1{,}45 I_z
\end{align*}
\begin{itemize}
\item $I_b$: corrente de projeto.
\item $I_n$: corrente nominal/ajuste do dispositivo.
\item $I_z$: capacidade do condutor após correções.
\item $I_2$: corrente convencional de atuação para sobrecarga.
\end{itemize}
\column{0.47\textwidth}
\centering
\begin{tikzpicture}[scale=0.9,every node/.style={font=\scriptsize}]
\draw[->] (0,0)--(5,0) node[right]{$I$};
\draw[thick,corDef] (1.0,0.1)--(1.0,1.1) node[above]{$I_b$};
\draw[thick,corAccent] (2.0,0.1)--(2.0,1.6) node[above]{$I_n$};
\draw[thick,corNorma] (3.3,0.1)--(3.3,1.2) node[above]{$I_z$};
\draw[decorate,decoration={brace,amplitude=5pt}] (1,-0.25)--(3.3,-0.25) node[midway,below=5pt]{coordenação};
\end{tikzpicture}
\end{columns}
\end{frame}

\begin{frame}{Fatores de correção térmica}
\small
\begin{tabularx}{\textwidth}{>{\bfseries}p{3.2cm}X}
\toprule
Fator & Consequência \\
\midrule
Temperatura ambiente & Aumenta temperatura de operação e reduz capacidade admissível. \\
Agrupamento & Circuitos próximos aquecem mutuamente. \\
Isolação térmica & Pode impedir dissipação e exigir forte redução de corrente. \\
Método de instalação & Eletroduto, leito, bandeja, enterrado e ao ar têm comportamentos diferentes. \\
Tipo de isolação & PVC, EPR/XLPE etc. possuem limites térmicos distintos. \\
Harmônicas & Aumentam perdas e podem elevar corrente no neutro. \\
\bottomrule
\end{tabularx}
\vspace{2mm}
\begin{caixaAt}
A corrente de catálogo do cabo não pode ser copiada sem identificar o método de referência e os fatores aplicáveis.
\end{caixaAt}
\end{frame}

\begin{frame}{Neutro em presença de harmônicas triplas}
\small
\begin{columns}[T,onlytextwidth]
\column{0.50\textwidth}
\begin{itemize}
\item Em sistema equilibrado senoidal, a fundamental do neutro tende a se cancelar.
\item Componentes de ordem 3, 9, 15... são de sequência zero e podem \textbf{somar} no neutro.
\item Fontes chaveadas monofásicas em grande quantidade exigem atenção especial.
\end{itemize}
\begin{caixaAt}
Não reduzir o neutro automaticamente em escritórios, data centers, laboratórios ou instalações com muita eletrônica.
\end{caixaAt}
\column{0.48\textwidth}
\begin{tikzpicture}
\begin{axis}[width=5.7cm,height=3.7cm,ybar,bar width=10pt,xlabel={ordem},ylabel={$I_h/I_1$},ymin=0,ymax=1.1,xtick={1,3,5,7,9},grid=major]
\addplot coordinates {(1,1)(3,0.55)(5,0.20)(7,0.14)(9,0.12)};
\end{axis}
\end{tikzpicture}
\end{columns}
\end{frame}

\begin{frame}[shrink=5]{Queda de tensão}
\small
\begin{columns}[T,onlytextwidth]
\column{0.53\textwidth}
\begin{align*}
\Delta V_{1\phi}&\approx2IL(R'\cos\varphi+X'\sen\varphi)\\
\Delta V_{3\phi}&\approx\sqrt3IL(R'\cos\varphi+X'\sen\varphi)\\
\Delta V(\%)&=100\frac{\Delta V}{V_n}
\end{align*}
\begin{itemize}
\item $L$ é o comprimento de ida; $R'$ e $X'$ são valores por unidade de comprimento, na condição de operação considerada. Usar unidades compatíveis, por exemplo km com $\Omega$/km.
\item Em $1\phi$, o fator 2 já representa ida e retorno; não duplicar novamente o comprimento.
\item Usar $V_n=V_{FN}$ no circuito monofásico e $V_n=V_{LL}$ no trifásico.
\item Circuitos longos podem ser dimensionados pela queda, não pela ampacidade.
\end{itemize}
\column{0.45\textwidth}
\begin{caixaEx}
Trifásico: $I=\SI{40}{A}$, $L=\SI{80}{m}$, $R'=\SI{3.08}{\ohm/km}$, $X'=\SI{0.08}{\ohm/km}$, $fp=0{,}85$.
\[
\Delta V\approx\SI{14.7}{V}\Rightarrow 3{,}9\%\;@\;\SI{380}{V}.
\]
A aceitação depende do limite aplicável ao trecho e ao sistema completo.
\end{caixaEx}
\end{columns}
\end{frame}

\begin{frame}{Efeito do comprimento: corrente de falta e queda}
\small
\centering
\begin{tikzpicture}
\begin{axis}[width=10.8cm,height=4.6cm,xlabel={distância ao quadro/fonte (m)},ylabel={valor normalizado},xmin=0,xmax=250,ymin=0,ymax=1.1,grid=both,legend pos=north east]
\addplot[corPrim,thick,smooth] coordinates {(0,1)(25,0.82)(50,0.67)(75,0.56)(100,0.47)(150,0.36)(200,0.29)(250,0.24)};
\addlegendentry{$I_{cc}$ relativo}
\addplot[corAccent,thick,dashed,smooth] coordinates {(0,0)(25,0.10)(50,0.20)(75,0.30)(100,0.40)(150,0.60)(200,0.80)(250,1.00)};
\addlegendentry{$\Delta V$ relativo}
\end{axis}
\end{tikzpicture}
\begin{caixaObs}
Quanto mais distante: maior queda de tensão e, em geral, menor corrente de curto. Isso pode alterar seção e tempo de desligamento.
\end{caixaObs}
\end{frame}

\begin{frame}{Curto-circuito: verificação térmica do condutor}
\small
\begin{columns}[T,onlytextwidth]
\column{0.51\textwidth}
\begin{align*}
S &\ge \frac{I_k\sqrt{t}}{k}
\end{align*}
\begin{itemize}
\item $S$: seção do condutor.
\item $I_k$: corrente de falta considerada.
\item $t$: duração até eliminação.
\item $k$: constante relacionada ao material e isolação.
\end{itemize}
\column{0.47\textwidth}
\begin{caixaEx}
Se um PE conduz $\SI{4}{kA}$ por $\SI{0.2}{s}$ e o valor adotado para $k$ for $115$:
\[
S\ge\frac{4000\sqrt{0{,}2}}{115}\approx\SI{15.6}{mm^2}.
\]
Escolhe-se seção comercial e verifica-se o conjunto de critérios.
\end{caixaEx}
\end{columns}
\end{frame}

\begin{frame}{Condutor de proteção: seção também é calculada}
\small
\begin{columns}[T,onlytextwidth]
\column{0.50\textwidth}
\begin{caixaNorma}[title=Método tabular simplificado]
Para condutores de fase e PE do mesmo material e nas condições previstas pelo método:
\[
S_{PE}=\begin{cases}
S_F, & S_F\le\SI{16}{mm^2}\\
\SI{16}{mm^2}, & \SI{16}{mm^2}<S_F\le\SI{35}{mm^2}\\
S_F/2, & S_F>\SI{35}{mm^2}
\end{cases}
\]
\end{caixaNorma}
\column{0.48\textwidth}
\begin{caixaAt}[title=O que ainda precisa fechar]
\begin{itemize}
\item seção mínima e proteção mecânica;
\item condição adiabática $S\ge I_k\sqrt{t}/k$;
\item impedância do caminho de falta e tempo de desligamento;
\item conexões, continuidade e identificação;
\item requisitos específicos de equipotencialização.
\end{itemize}
\end{caixaAt}
\end{columns}
\end{frame}

\begin{frame}{Eletrodutos, bandejas e segregação}
\small
\begin{columns}[T,onlytextwidth]
\column{0.50\textwidth}
\begin{itemize}
\item Ocupação deve permitir lançamento e substituição de cabos.
\item Respeitar raios de curvatura e esforços de tração.
\item Separar potência, controle, instrumentação e dados conforme compatibilidade e isolamento.
\item Considerar expansão e acessibilidade.
\item Bandejas metálicas precisam integrar o conceito de equipotencialização quando aplicável.
\end{itemize}
\column{0.48\textwidth}
\centering
\begin{tikzpicture}[scale=0.78]
\draw[thick,corPrim] (0,0) rectangle (5.8,2.1);
\foreach \x in {0.25,0.65,...,5.55}{\draw[black!25] (\x,0)--(\x,2.1);}
\foreach \x/\y/\c in {0.7/1.45/corPrim,1.5/0.7/corAccent,2.5/1.3/corAt,3.4/0.55/corNorma,4.3/1.5/corDef,5.1/0.85/corPrim}{\fill[\c!70] (\x,\y) circle (0.18);\draw (\x,\y) circle (0.18);}
\node at (2.9,-0.45) {leito/bandeja: rota também é parte do projeto};
\end{tikzpicture}
\end{columns}
\end{frame}

\begin{frame}{Identificação de condutores e circuitos}
\small
\begin{columns}[T,onlytextwidth]
\column{0.50\textwidth}
\begin{tabularx}{\textwidth}{>{\bfseries}p{2.1cm}X}
\toprule
Elemento & Princípio \\
\midrule
PE & Identificação exclusiva de proteção. \\
Neutro & Identificação própria e coerente. \\
Fases & Padrão consistente em toda a instalação. \\
Retornos & Anilhas/etiquetas; não confundir com PE/N. \\
Cabos & Origem, destino, circuito e seção. \\
\bottomrule
\end{tabularx}
\column{0.48\textwidth}
\begin{caixaProjeto}[title=As built útil]
Registrar:
\begin{itemize}
\item identificação real;
\item bornes e dispositivos;
\item ajustes de proteção;
\item cabos instalados;
\item alterações de campo;
\item reservas e circuitos desativados.
\end{itemize}
\end{caixaProjeto}
\end{columns}
\end{frame}

%=====================================================================

\section{Proteção, DR, DPS e aterramento}

\begin{frame}{Proteções: cada dispositivo resolve um problema}
\scriptsize
\begin{tabularx}{\textwidth}{>{\bfseries}p{3.0cm}X}
\toprule
Dispositivo & Função principal \\
\midrule
Disjuntor/fusível & Sobrecarga e curto-circuito; deve ter capacidade de interrupção adequada. \\
DR/IDR/RCBO & Corrente diferencial residual: choque/fuga conforme função, tipo e sensibilidade. \\
DPS & Limita sobretensões transitórias; atua em coordenação com aterramento, equipotencialização e SPDA. \\
Seccionadora & Isolação/manobra para operação e manutenção conforme aplicação. \\
Contator & Manobra frequente; não substitui proteção de curto. \\
Relé de motor & Sobrecarga, falta de fase, travamento e outras funções. \\
Proteção eletrônica & Ajustes de longo/curto/instantâneo, falta à terra, comunicação etc. \\
\bottomrule
\end{tabularx}
\end{frame}

\begin{frame}{Proteção contra choque: camadas}
\small
\begin{columns}[T,onlytextwidth]
\column{0.31\textwidth}
\begin{caixaDef}[title=Proteção básica]
Evitar contato com partes vivas em condição normal: isolação, barreiras, invólucros etc.
\end{caixaDef}
\column{0.31\textwidth}
\begin{caixaNorma}[title=Proteção supletiva]
Limitar risco se ocorrer falha: PE, equipotencialização, desligamento automático etc.
\end{caixaNorma}
\column{0.34\textwidth}
\begin{caixaAt}[title=Proteção adicional]
DR de alta sensibilidade em situações previstas. É complemento, não substituto das medidas básicas.
\end{caixaAt}
\end{columns}
\vspace{3mm}
\begin{caixaObs}
A proteção contra choque é um sistema de medidas coordenadas; não existe ``resolver tudo instalando um DR''.
\end{caixaObs}
\end{frame}

\begin{frame}{DR: princípio de funcionamento}
\small
\begin{columns}[T,onlytextwidth]
\column{0.50\textwidth}
\begin{caixaDef}
O DR compara a soma vetorial das correntes dos condutores ativos que atravessam seu toroide. Uma fuga que não retorna por esses condutores produz corrente residual.
\end{caixaDef}
\[
I_{\Delta}=\left|\sum I_k\right|.
\]
\begin{itemize}
\item Neutros de circuitos protegidos não podem ser misturados a jusante.
\item PE não atravessa o toroide como condutor ativo.
\end{itemize}
\column{0.48\textwidth}
\begin{circuitikz}[scale=0.76,every node/.style={font=\scriptsize}]
\draw (0,3) node[left]{F} -- (1,3) to[generic,l=DR] (2.4,3) -- (4.8,3) to[R,l=Carga] (4.8,1.2) -- (2.4,1.2) to[generic] (1,1.2) -- (0,1.2) node[left]{N};
\draw[->,thick,corPrim] (2.6,3.2)--(3.6,3.2) node[midway,above]{$I_F$};
\draw[->,thick,corPrim] (3.6,1.0)--(2.6,1.0) node[midway,below]{$I_N$};
\draw[->,thick,corAt] (4.8,2.1)--(5.6,1.2) node[right]{fuga};
\draw[corNorma,thick] (5.6,1.2)--(5.6,0.55)--(0.6,0.55) node[left]{PE};
\end{circuitikz}
\end{columns}
\end{frame}

\begin{frame}[shrink=8]{Tipos de DR e eletrônica de potência}
\scriptsize
\begin{tabularx}{\textwidth}{>{\bfseries}p{1.4cm}p{4.3cm}X}
\toprule
Tipo & Sensibilidade funcional & Aplicações/alertas \\
\midrule
AC & Corrente residual senoidal CA. & Cargas simples; pode ser inadequado para eletrônica moderna. \\
A & CA senoidal + componentes pulsantes CC previstas. & Eletrodomésticos e muitas cargas eletrônicas. \\
F & Inclui características para determinadas cargas monofásicas com conversores de frequência. & Alguns equipamentos inverter/velocidade variável. \\
B & Detecta também correntes residuais CC lisas dentro do escopo do dispositivo. & Pode ser requerido em VFD, FV, VE e outros conversores; não é escolha automática sem analisar o equipamento e a norma específica. \\
\bottomrule
\end{tabularx}
\vspace{2mm}
\begin{caixaAt}
A escolha não deve ser feita apenas por ``30 mA''. Verificar \textbf{tipo de forma de onda residual}, seletividade, imunidade, corrente de fuga normal, detecção de CC incorporada ao equipamento e norma específica. Em VE, por exemplo, uma solução com detecção de \SI{6}{mA} CC incorporada não é equivalente a instalar indiscriminadamente qualquer DR tipo A.
\end{caixaAt}
\end{frame}

\begin{frame}[shrink=4]{DPS: parâmetros e topologia que importam}
\small
\begin{columns}[T,onlytextwidth]
\column{0.47\textwidth}
\begin{itemize}
\item $U_c$: máxima tensão contínua aplicável.
\item $U_p$: nível de proteção de tensão.
\item $I_n$: corrente nominal de descarga em ensaio aplicável.
\item $I_{imp}$: corrente de impulso para DPS tipo 1.
\item Esquema de conexão compatível com o aterramento da instalação.
\item Proteção de retaguarda e capacidade de curto conforme fabricante.
\item Condutores de conexão tão curtos e retilíneos quanto praticável.
\end{itemize}
\column{0.51\textwidth}
\centering
\begin{tikzpicture}[scale=0.78,every node/.style={font=\scriptsize}]
\draw[thick] (0,3.0)--(5.3,3.0); \node[left] at (0,3.0){F};
\draw[thick] (0,2.1)--(5.3,2.1); \node[left] at (0,2.1){N};
\draw[thick,corNorma] (0,0.7)--(5.3,0.7); \node[left,corNorma] at (0,0.7){PE};
\node[draw=corAt,fill=corAt!8,rounded corners,minimum width=1.05cm,minimum height=0.55cm] (d1) at (2.0,1.55){DPS};
\node[draw=corAt,fill=corAt!8,rounded corners,minimum width=1.05cm,minimum height=0.55cm] (d2) at (3.7,1.55){DPS};
\draw[corAt,thick] (2.0,3.0)--(d1.north); \draw[corNorma,thick] (d1.south)--(2.0,0.7);
\draw[corAt,thick] (3.7,2.1)--(d2.north); \draw[corNorma,thick] (d2.south)--(3.7,0.7);
\node[font=\tiny,align=center] at (2.85,0.15){PE segue ao BEP / sistema\\de equipotencialização};
\end{tikzpicture}
\end{columns}
\begin{caixaObs}
O desenho é conceitual. A combinação F--PE, N--PE, F--N ou 3+1 depende do esquema de aterramento, do DPS selecionado e da coordenação requerida. Não se cria um ``terra local'' isolado para o DPS.
\end{caixaObs}
\end{frame}

\begin{frame}{Coordenação de DPS em cascata}
\small
\centering
\begin{tikzpicture}[scale=0.83,every node/.style={font=\scriptsize}]
\tikzset{b/.style={draw=corPrim,fill=corDef!7,rounded corners,minimum width=2.5cm,minimum height=0.7cm,align=center}}
\node[b] (a) at (0,2.2) {Entrada / QGBT\\DPS tipo 1 ou 1+2};
\node[b] (b) at (4.0,2.2) {QD\\DPS tipo 2};
\node[b] (c) at (8.0,2.2) {Equipamento\\proteção fina};
\draw[-{Stealth},thick,corPrim] (a)--(b);
\draw[-{Stealth},thick,corPrim] (b)--(c);
\node[font=\tiny\bfseries,corPrim,fill=white,inner sep=1pt] at (2.0,2.82){alimentador};
\node[font=\tiny\bfseries,corPrim,fill=white,inner sep=1pt] at (6.0,2.82){circuito};
% PE/equipotencialização contínuos, sem eletrodos independentes
\draw[corNorma,very thick] (-0.7,0.7)--(8.7,0.7);
\node[left,corNorma] at (-0.7,0.7){PE/BEP};
\draw[corNorma,thick] (0,1.85)--(0,0.7);
\draw[corNorma,thick] (4,1.85)--(4,0.7);
\draw[corNorma,thick] (8,1.85)--(8,0.7);
\draw[corNorma,thick] (1.0,0.7)--(1.0,-0.1) node[ground]{};
\node[font=\tiny,corNorma] at (2.35,0.25){sistema comum de equipotencialização/aterramento};
\end{tikzpicture}
\vspace{2mm}
\begin{caixaNorma}[title=Coordenação]
A cascata deve considerar categoria/local da proteção, tensão suportável do equipamento, distância, energia do surto, presença de SPDA, esquema de aterramento e instruções dos fabricantes.
\end{caixaNorma}
\end{frame}

\begin{frame}[shrink=2]{Funções do aterramento}
\small
\begin{columns}[T,onlytextwidth]
\column{0.52\textwidth}
\begin{itemize}
\item Referência e caminho de corrente de falta conforme esquema.
\item Limitação de tensões de toque e passo.
\item Viabilização do desligamento automático.
\item Equipotencialização de massas e elementos condutivos.
\item Dissipação/controle de surtos em conjunto com DPS/SPDA.
\item Compatibilidade eletromagnética e referência funcional quando aplicável.
\end{itemize}
\column{0.46\textwidth}
\centering
\begin{tikzpicture}[scale=0.82,every node/.style={font=\scriptsize}]
\draw[thick] (0,0) rectangle (3.8,2.0);
\node at (1.9,1.0) {Quadro / BEP};
\draw[thick,corNorma] (1.9,0) -- (1.9,-0.8) -- (0.6,-0.8) -- (0.6,-1.8);
\draw[thick,corNorma] (1.9,-0.8) -- (3.2,-0.8) -- (3.2,-1.8);
\foreach \x in {0.6,3.2}{\draw[thick,corNorma] (\x,-1.8)--(\x,-2.7);\draw[thick,corNorma] (\x-0.3,-2.7)--(\x+0.3,-2.7);}
\draw[corNorma,thick] (0.6,-2.25)--(3.2,-2.25);
\node[corNorma] at (1.9,-3.05) {eletrodos / malha};
\end{tikzpicture}
\end{columns}
\begin{caixaAt}
``Aterramento = uma haste'' é uma simplificação perigosa. O desempenho depende do esquema completo, impedâncias, equipotencialização e dispositivos de proteção.
\end{caixaAt}
\end{frame}

\begin{frame}{Esquemas TN, TT e IT}
\scriptsize
\begin{tabularx}{\textwidth}{>{\bfseries}p{1.5cm}p{3.5cm}X}
\toprule
Esquema & Característica & Pontos de projeto \\
\midrule
TN-S & N e PE separados. & Boa previsibilidade da corrente de falta; atenção a laços e EMC. \\
TN-C & PEN combinado em parte permitida do sistema. & Restrições importantes; não é solução genérica para circuitos terminais. \\
TN-C-S & PEN a montante e separação posterior de N/PE. & Ponto de separação deve ser controlado e não se recombina N com PE a jusante. \\
TT & Massas com eletrodo local, distinto funcionalmente do aterramento da fonte. & DR costuma ser essencial para desligamento automático. \\
IT & Fonte isolada da terra ou ligada por impedância. & Alta continuidade na primeira falta; exige monitoramento de isolamento e estratégia de localização. \\
\bottomrule
\end{tabularx}
\end{frame}

\begin{frame}{Equipotencialização principal e suplementar}
\small
\begin{columns}[T,onlytextwidth]
\column{0.49\textwidth}
\begin{caixaDef}
Equipotencializar é reduzir diferenças perigosas de potencial interligando, de forma planejada, massas e elementos condutivos relevantes ao sistema de equipotencialização.
\end{caixaDef}
\begin{itemize}
\item BEP na origem apropriada.
\item PE e eletrodo(s) de aterramento.
\item Partes metálicas relevantes e serviços que ingressem na edificação.
\item Ligações suplementares onde necessárias.
\end{itemize}
\column{0.49\textwidth}
\centering
\begin{tikzpicture}[scale=0.78,every node/.style={font=\scriptsize}]
\node[draw,fill=corDef!8] (bep) at (0,2) {BEP};
\node[draw] (qd) at (4,2) {QD};
\node[draw] (tub) at (0,0) {Tubulação};
\node[draw] (maq) at (4,0) {Massa};
\node[draw] (est) at (2,-1.2) {Estrutura};
\foreach \x in {qd,tub,maq,est}{\draw[thick,corNorma] (bep)--(\x);}
\draw[thick,corNorma] (bep)--(-1,1) node[ground]{};
\end{tikzpicture}
\end{columns}
\end{frame}

%=====================================================================

\section{SPDA e proteção contra descargas atmosféricas}

\begin{frame}{NBR 5419: atualização confirmada em 2026}
\small
\begin{columns}[T,onlytextwidth]
\column{0.52\textwidth}
\begin{caixaNorma}[title=ABNT NBR 5419-1:2026]
A edição 2026 da \textbf{Parte 1} está publicada e estabelece os requisitos para a determinação da necessidade de proteção contra descargas atmosféricas.
\end{caixaNorma}
\begin{itemize}
\item não assumir automaticamente que todas as partes da série mudaram de edição na mesma data;
\item no projeto, registrar a edição efetivamente adotada de cada parte aplicável;
\item verificar requisitos do órgão licenciador, Corpo de Bombeiros e contratante.
\end{itemize}
\column{0.46\textwidth}
\begin{caixaAt}[title=Regra para este material]
A referência ``NBR 5419:2015'' não é mais usada de forma genérica. Quando a Parte 1 for citada, o material indica explicitamente \textbf{ABNT NBR 5419-1:2026}; para as demais partes, exige-se conferência da edição vigente antes do projeto executivo.
\end{caixaAt}
\end{columns}
\end{frame}

\begin{frame}{PDA não é apenas o captor no telhado}
\small
\centering
\begin{tikzpicture}[scale=0.82,every node/.style={font=\scriptsize}]
\tikzset{b/.style={draw=corPrim,fill=corDef!7,rounded corners,minimum width=3.0cm,minimum height=0.7cm,align=center}}
\node[b] (r) at (0,2.6) {Avaliação da necessidade\\e análise de risco};
\node[b] (s) at (-3.4,1.1) {SPDA externo\\captação/descida/terra};
\node[b] (e) at (0,1.1) {Equipotencialização\\para descargas};
\node[b] (m) at (3.4,1.1) {MPS\\DPS, roteamento, blindagem};
\node[b] (i) at (0,-0.5) {Sistemas internos\\energia, dados e automação};
\draw[-{Stealth},thick,corPrim] (r)--(s);\draw[-{Stealth},thick,corPrim] (r)--(e);\draw[-{Stealth},thick,corPrim] (r)--(m);
\draw[-{Stealth},thick,corNorma] (s)--(i);\draw[-{Stealth},thick,corNorma] (e)--(i);\draw[-{Stealth},thick,corNorma] (m)--(i);
\end{tikzpicture}
\vspace{2mm}
\begin{caixaObs}
A interface entre SPDA, equipotencialização e DPS precisa ser projetada conjuntamente. Separar esses projetos em ``ilhas'' cria lacunas de proteção.
\end{caixaObs}
\end{frame}

\begin{frame}{Zonas de proteção e roteamento}
\small
\begin{columns}[T,onlytextwidth]
\column{0.49\textwidth}
\begin{itemize}
\item A energia do surto deve ser reduzida progressivamente entre zonas.
\item Entradas metálicas e cabos são pontos críticos de acoplamento.
\item Roteamento, blindagem, equipotencialização e DPS trabalham em conjunto.
\item Cabos sensíveis não devem percorrer rotas inadequadas junto a condutores de descida sem análise.
\end{itemize}
\column{0.49\textwidth}
\centering
\begin{tikzpicture}[scale=0.82,every node/.style={font=\scriptsize}]
\draw[thick,corAt] (0,0) rectangle (5.3,3.4);
\draw[thick,corAccent] (0.7,0.6) rectangle (4.6,2.8);
\draw[thick,corNorma] (1.5,1.2) rectangle (3.8,2.2);
\node[corAt] at (4.65,3.1) {ZPR externa};
\node[corAccent] at (3.9,2.55) {zona interna};
\node[corNorma] at (2.65,1.7) {equipamento sensível};
\draw[-{Stealth},thick,corPrim] (-0.6,1.7)--(1.4,1.7) node[midway,above]{linha};
\node at (0.7,1.35) {DPS};
\end{tikzpicture}
\end{columns}
\end{frame}

%=====================================================================

% Núcleo industrial
\section{Instalações industriais, QGBT, CCM e motores}

\begin{frame}{Arquitetura industrial típica}
\small
\centering
\begin{tikzpicture}[scale=0.83,every node/.style={font=\scriptsize}]
\tikzset{b/.style={draw=corPrim,fill=corDef!7,rounded corners,minimum width=2.45cm,minimum height=0.67cm,align=center}}
\node[b] (mt) at (0,3.2) {Rede MT};
\node[b] (sub) at (0,2.2) {Subestação / trafo};
\node[b] (qgbt) at (0,1.2) {QGBT};
\node[b] (ccm) at (-3.4,0.0) {CCM / motores};
\node[b] (qd) at (0,0.0) {QD / processo};
\node[b] (ups) at (3.4,0.0) {UPS / críticas};
\node[b] (m) at (-3.4,-1.2) {Bombas / ventiladores};
\node[b] (p) at (0,-1.2) {Força / iluminação};
\node[b] (clp) at (3.4,-1.2) {CLP / TI / instrumentos};
\foreach \u/\v in {mt/sub,sub/qgbt,qgbt/ccm,qgbt/qd,qgbt/ups,ccm/m,qd/p,ups/clp}{\draw[-{Stealth},thick,corPrim] (\u)--(\v);}
\end{tikzpicture}
\vspace{2mm}
\begin{caixaObs}
Na indústria, a decisão técnica precisa incluir disponibilidade, manutenção, expansão e custo da parada de processo.
\end{caixaObs}
\end{frame}

\begin{frame}{QGBT: especificação técnica}
\small
\begin{columns}[T,onlytextwidth]
\column{0.52\textwidth}
\begin{itemize}
\item Tensão e corrente nominal do conjunto.
\item Corrente de curto suportável e barramentos.
\item Grau de proteção e condições ambientais.
\item Segregação/compartimentação conforme necessidade.
\item Dissipação térmica e ventilação.
\item Medição, supervisão, comunicação e reservas.
\end{itemize}
\column{0.46\textwidth}
\begin{caixaNorma}[title=Série NBR IEC 61439]
O quadro deve ser tratado como \textbf{conjunto}, com responsabilidades de projeto/montagem e verificações aplicáveis -- não como simples reunião de dispositivos individuais certificados.
\end{caixaNorma}
\begin{caixaAt}
Icc do dispositivo não substitui a análise da suportabilidade do conjunto e dos barramentos.
\end{caixaAt}
\end{columns}
\end{frame}

\begin{frame}{CCM: centro de controle de motores}
\small
\begin{columns}[T,onlytextwidth]
\column{0.49\textwidth}
\begin{itemize}
\item Alimentação e distribuição para motores.
\item Proteção, seccionamento e manobra.
\item Partida direta, soft-starter, inversor.
\item Relés eletrônicos e monitoramento.
\item Interface com CLP, IHM e supervisório.
\item Intertravamentos e segurança operacional.
\end{itemize}
\column{0.49\textwidth}
\centering
\begin{tikzpicture}[scale=0.72,every node/.style={font=\tiny}]
\draw[thick,fill=black!4] (0,0) rectangle (6.0,3.8);
\foreach \x in {1.5,3.0,4.5}{\draw (\x,0)--(\x,3.8);}
\foreach \y in {0.75,1.5,2.25,3.0}{\draw (0,\y)--(6.0,\y);}
\foreach \x in {0.75,2.25,3.75,5.25}{\foreach \y in {0.38,1.13,1.88,2.63,3.38}{\node at (\x,\y){M};}}
\node[font=\scriptsize] at (3.0,-0.4) {unidades funcionais / gavetas};
\end{tikzpicture}
\end{columns}
\end{frame}

\begin{frame}{Motor de indução: corrente e partida}
\small
\begin{columns}[T,onlytextwidth]
\column{0.51\textwidth}
\begin{align*}
I_N&=\frac{P_{out}}{\sqrt3V\eta fp}\\
I_{part}&\approx k_p I_N
\end{align*}
\begin{itemize}
\item $k_p$ depende do motor e método de partida.
\item Partida afeta queda de tensão e dimensionamento da rede.
\item Proteção precisa considerar sobrecarga, curto, falta de fase e travamento.
\end{itemize}
\column{0.47\textwidth}
\begin{tikzpicture}
\begin{axis}[width=5.7cm,height=3.9cm,xlabel={tempo},ylabel={$I/I_N$},xmin=0,xmax=6,ymin=0,ymax=7,grid=both,legend pos=north east]
\addplot[corPrim,thick] coordinates {(0,6.2)(0.5,5.8)(1,4.5)(2,2.2)(3,1.2)(6,1)};\addlegendentry{direta}
\addplot[corAccent,thick,dashed] coordinates {(0,2.8)(0.5,2.5)(1,2)(2,1.4)(3,1.05)(6,1)};\addlegendentry{rampa}
\end{axis}
\end{tikzpicture}
\end{columns}
\end{frame}

\begin{frame}{Comparação de métodos de partida}
\scriptsize
\begin{tabularx}{\textwidth}{>{\bfseries}p{2.5cm}p{2.7cm}p{2.8cm}X}
\toprule
Método & Corrente de partida & Torque/controle & Quando faz sentido \\
\midrule
Direta & Alta & Torque natural do motor & Rede robusta, motor pequeno/médio e processo tolerante. \\
Estrela-triângulo & Reduzida & Torque também reduzido & Motor compatível e carga que acelera com torque reduzido. \\
Soft-starter & Rampa de tensão & Partida/parada suave & Bombas, ventiladores e redução de impacto mecânico/hidráulico. \\
Inversor & Controlada & Velocidade e torque variáveis & Processo exige controle, economia e flexibilidade. \\
\bottomrule
\end{tabularx}
\vspace{2mm}
\begin{caixaAt}
Não escolher pelo preço do equipamento isoladamente: comparar rede, carga mecânica, ciclo, eficiência, harmônicas, manutenção e controle.
\end{caixaAt}
\end{frame}

\begin{frame}{Inversor de frequência: efeitos na instalação}
\small
\begin{columns}[T,onlytextwidth]
\column{0.51\textwidth}
\begin{itemize}
\item Retificador gera correntes não senoidais na entrada.
\item PWM na saída cria alto $dv/dt$ e correntes de modo comum.
\item Comprimento do cabo ao motor pode exigir filtro/reator.
\item Blindagem e aterramento afetam EMC.
\item Correntes de fuga interferem na escolha do DR.
\end{itemize}
\column{0.47\textwidth}
\centering
\begin{tikzpicture}[every node/.style={font=\scriptsize}]
\tikzset{b/.style={draw=corPrim,fill=corDef!7,rounded corners,minimum width=2.2cm,minimum height=0.58cm}}
\node[b] (rede) at (0,4.4) {Rede CA};
\node[b] (ret) at (0,3.3) {Retificador};
\node[b] (dc) at (0,2.2) {Link CC};
\node[b] (inv) at (0,1.1) {PWM};
\node[b] (mot) at (0,0) {Motor};
\foreach \u/\v in {rede/ret,ret/dc,dc/inv,inv/mot}{\draw[-{Stealth},thick,corPrim] (\u)--(\v);}
\draw[corNorma,thick] (mot.south)--(0,-0.8) node[ground]{};
\end{tikzpicture}
\end{columns}
\end{frame}

%=====================================================================

\section{Curto-circuito, seletividade e arco elétrico}

\begin{frame}{Corrente de curto-circuito presumida}
\small
\begin{columns}[T,onlytextwidth]
\column{0.50\textwidth}
\begin{align*}
I_{cc,3\phi}&\approx\frac{V_{LL}}{\sqrt3 Z_{eq}}\\
I_{cc,trafo}&\approx\frac{I_N}{Z_{pu}}
\end{align*}
\begin{itemize}
\item Fonte, transformador, cabos e conexões entram em $Z_{eq}$.
\item Calcular curto máximo e mínimo conforme finalidade.
\end{itemize}
\column{0.48\textwidth}
\begin{caixaEx}
Trafo $\SI{300}{kVA}$, $\SI{380}{V}$, $Z=5\%$:
\[
I_N=\frac{300000}{\sqrt3\cdot380}=\SI{456}{A}
\]
\[
I_{cc}\approx\frac{456}{0{,}05}=\SI{9.1}{kA}.
\]
É uma aproximação no secundário do trafo, antes da impedância dos cabos.
\end{caixaEx}
\end{columns}
\end{frame}

\begin{frame}{Capacidade de interrupção e suportabilidade}
\small
\begin{columns}[T,onlytextwidth]
\column{0.49\textwidth}
\begin{caixaAt}[title=Dispositivo]
A capacidade de interrupção do disjuntor/fusível deve ser compatível com a corrente de curto presumida no ponto.
\end{caixaAt}
\column{0.49\textwidth}
\begin{caixaNorma}[title=Conjunto]
Barramentos, cabos e quadros precisam suportar os efeitos térmicos/dinâmicos até a eliminação da falta.
\end{caixaNorma}
\end{columns}
\vspace{3mm}
\begin{caixaObs}
``Disjuntor 10 kA'' não garante que todo QD/QGBT tenha suportabilidade de 10 kA. O conjunto deve ser especificado e verificado como sistema.
\end{caixaObs}
\end{frame}

\begin{frame}{Seletividade: manter apenas a parte defeituosa fora}
\small
\begin{columns}[T,onlytextwidth]
\column{0.49\textwidth}
\begin{itemize}
\item Seletividade amperimétrica.
\item Seletividade temporal.
\item Seletividade energética.
\item Zonas/intertravamentos em proteção eletrônica.
\item Tabelas de coordenação do fabricante quando aplicáveis.
\end{itemize}
\column{0.49\textwidth}
\centering
\begin{tikzpicture}[scale=0.82,every node/.style={font=\scriptsize}]
\tikzset{d/.style={draw=corPrim,fill=corDef!7,rounded corners,minimum width=1.7cm,minimum height=0.55cm}}
\node[d] (dg) at (0,3) {DJ geral};
\node[d] (d1) at (-2.2,1.8) {DJ QD1};
\node[d] (d2) at (2.2,1.8) {DJ QD2};
\node[d] (c1) at (-2.2,0.6) {Carga 1};
\node[d] (c2) at (2.2,0.6) {Carga 2};
\draw[-{Stealth},thick] (dg)--(d1);\draw[-{Stealth},thick] (dg)--(d2);\draw[-{Stealth},thick] (d1)--(c1);\draw[-{Stealth},thick] (d2)--(c2);
\draw[corAt,thick,decorate,decoration={zigzag,segment length=4pt,amplitude=2pt}] (c1.south) -- ++(0,-0.65) node[below]{falta};
\node[corAt,align=center] at (0,0.15) {atua\\QD1};
\end{tikzpicture}
\end{columns}
\end{frame}

\begin{frame}{Curvas tempo--corrente}
\small
\begin{columns}[T,onlytextwidth]
\column{0.43\textwidth}
\begin{itemize}
\item Eixo horizontal: corrente ou múltiplo de $I_n$.
\item Eixo vertical: tempo de atuação.
\item Para seletividade, a curva a montante deve ficar adequadamente coordenada com a jusante no domínio de interesse.
\item Curto mínimo também importa: precisa haver atuação suficientemente rápida.
\end{itemize}
\column{0.55\textwidth}
\begin{tikzpicture}
\begin{axis}[width=6.3cm,height=4.2cm,xmode=log,ymode=log,xlabel={$I$},ylabel={$t$},xmin=1,xmax=100,ymin=0.01,ymax=1000,grid=both,legend pos=south west]
\addplot[corPrim,thick] coordinates {(1.1,800)(1.5,120)(2,25)(4,2)(8,0.05)};\addlegendentry{jusante}
\addplot[corAccent,thick,dashed] coordinates {(1.3,900)(2,180)(4,35)(8,3)(20,0.08)};\addlegendentry{montante}
\end{axis}
\end{tikzpicture}
\end{columns}
\end{frame}

\begin{frame}{Arco elétrico e energia incidente}
\small
\begin{columns}[T,onlytextwidth]
\column{0.52\textwidth}
\begin{caixaDef}
Arco elétrico é uma descarga sustentada através do ar/plasma. Pode gerar calor extremo, pressão, metal vaporizado, luz intensa e ruído.
\end{caixaDef}
\begin{itemize}
\item Risco não é avaliado apenas pela tensão nominal.
\item Tempo de atuação da proteção tem grande influência.
\item Distância de trabalho e configuração do equipamento importam.
\end{itemize}
\column{0.46\textwidth}
\begin{caixaNorma}[title=NR-10 futura --- vigência em 01/06/2027]
A redação aprovada pela Portaria MTE nº 737/2026, com vigência a partir de 01/06/2027, prevê, quando aplicável, que o projeto considere \textbf{estudo de energia incidente} para definir medidas de proteção contra os efeitos térmicos do arco elétrico. Até 31/05/2027, permanece vigente a redação anterior da NR-10.
\end{caixaNorma}
\end{columns}
\end{frame}

\begin{frame}{Como reduzir risco de arco elétrico}
\small
\begin{columns}[T,onlytextwidth]
\column{0.49\textwidth}
\begin{itemize}
\item Desenergizar sempre que possível.
\item Reduzir tempo de eliminação da falta.
\item Ajustes temporários de manutenção quando tecnicamente previstos.
\item Relés rápidos, ZSI, diferencial de barras, detecção de arco.
\item Operação remota e distância.
\end{itemize}
\column{0.49\textwidth}
\begin{itemize}
\item Compartimentação adequada.
\item Limitação de corrente quando aplicável.
\item Intertravamentos e portas adequadas.
\item Sinalização, análise de risco e procedimentos.
\item EPI baseado na avaliação e não como primeira medida de controle.
\end{itemize}
\end{columns}
\begin{caixaAt}
EPI é a última barreira. O projeto deve reduzir a energia e a probabilidade de exposição antes de depender do trabalhador.
\end{caixaAt}
\end{frame}

%=====================================================================


% Segurança e desempenho comuns
\section{Segurança, NR-10 e manutenção}

\begin{frame}{Desenergização: lógica de segurança}
\small
\centering
\begin{tikzpicture}[node distance=0.30cm,every node/.style={font=\scriptsize}]
\tikzset{b/.style={draw=corPrim,fill=corDef!7,rounded corners,minimum width=2.4cm,minimum height=0.65cm,align=center}}
\node[b] (a) {Seccionar};
\node[b,right=of a] (b) {Impedir\\reenergização};
\node[b,right=of b] (c) {Constatar ausência\\de tensão};
\node[b,below=of c] (d) {Aterramento\\temporário, se aplicável};
\node[b,left=of d] (e) {Proteger contra\\partes energizadas};
\node[b,left=of e] (f) {Sinalizar o impedimento\\de reenergização};
\foreach \u/\v in {a/b,b/c,c/d,d/e,e/f}{\draw[-{Stealth},thick,corPrim] (\u)--(\v);}
\end{tikzpicture}
\vspace{3mm}
\begin{caixaObs}
A sequência real deve seguir procedimento aprovado, análise de risco e requisitos aplicáveis. O slide mostra a lógica, não substitui procedimento operacional.
\end{caixaObs}
\end{frame}

\begin{frame}{LOTO: bloqueio e etiquetagem}
\small
\begin{columns}[T,onlytextwidth]
\column{0.50\textwidth}
\begin{itemize}
\item Identificar todas as fontes de energia.
\item Seccionar e bloquear cada fonte aplicável.
\item Dissipar/neutralizar energias armazenadas.
\item Etiquetar responsável e condição.
\item Verificar estado seguro antes do serviço.
\item Restabelecer somente após procedimento de liberação.
\end{itemize}
\column{0.48\textwidth}
\centering
\begin{tikzpicture}[scale=0.86,every node/.style={font=\scriptsize}]
\draw[thick,fill=corAt!8,rounded corners=3pt] (0,0) rectangle (4.8,3.2);
\node[font=\bfseries\large,corAt] at (2.4,2.65) {BLOQUEADO};
\draw[thick] (1.4,0.7) rectangle (3.4,1.8);
\draw[thick] (2.4,1.8) arc (180:0:0.7 and 0.85);
\node at (2.4,0.35) {não operar};
\end{tikzpicture}
\end{columns}
\end{frame}

\begin{frame}{Prontuário e documentação de segurança}
\small
\begin{columns}[T,onlytextwidth]
\column{0.49\textwidth}
\begin{itemize}
\item Diagramas atualizados.
\item Procedimentos e instruções técnicas.
\item Inspeções e medições do sistema de proteção.
\item Registros de treinamento/autorização conforme aplicável.
\item Especificação de EPC/EPI.
\end{itemize}
\column{0.49\textwidth}
\begin{itemize}
\item Estudos de risco pertinentes.
\item Relatórios de ensaio e manutenção.
\item Histórico de modificações.
\item Dados de curto-circuito e ajustes quando necessários à segurança.
\item Plano de emergência aplicável.
\end{itemize}
\end{columns}
\end{frame}

\begin{frame}{Termografia: exemplo de interpretação}
\small
\begin{columns}[T,onlytextwidth]
\column{0.48\textwidth}
\begin{itemize}
\item Comparar componentes equivalentes sob carga semelhante.
\item Investigar corrente, aperto, oxidação, dimensionamento e ventilação.
\item Temperatura absoluta isolada não explica a causa.
\item Registrar carga, ambiente e condições da inspeção.
\end{itemize}
\column{0.50\textwidth}
\begin{tikzpicture}
\begin{axis}[ybar,width=5.8cm,height=3.8cm,bar width=16pt,ylabel={$T$ (°C)},symbolic x coords={A,B,C,N},xtick=data,ymin=20,ymax=100,nodes near coords,grid=major]
\addplot coordinates {(A,48)(B,51)(C,82)(N,55)};
\end{axis}
\end{tikzpicture}
\end{columns}
\begin{caixaAt}
Fase C muito mais quente é um \textbf{indício}; a ação correta é diagnosticar antes de apenas ``reapertar tudo''.
\end{caixaAt}
\end{frame}

%=====================================================================

\section{Qualidade de energia e eficiência}

\begin{frame}{Harmônicas: medida básica}
\small
\begin{columns}[T,onlytextwidth]
\column{0.50\textwidth}
\[
THD_I=100\frac{\sqrt{I_2^2+I_3^2+\cdots+I_n^2}}{I_1}.
\]
\begin{itemize}
\item THD não informa sozinho qual harmônica domina.
\item Analisar espectro e forma de onda.
\item Fonte, carga e impedância da rede determinam a distorção de tensão.
\end{itemize}
\column{0.48\textwidth}
\begin{tikzpicture}
\begin{axis}[width=5.8cm,height=3.8cm,ybar,bar width=10pt,xlabel={ordem},ylabel={$I_h/I_1$},ymin=0,ymax=1.1,xtick={1,3,5,7,9,11},grid=major]
\addplot coordinates {(1,1)(3,0.42)(5,0.22)(7,0.16)(9,0.09)(11,0.07)};
\end{axis}
\end{tikzpicture}
\end{columns}
\end{frame}

\begin{frame}{Efeitos de harmônicas}
\small
\begin{columns}[T,onlytextwidth]
\column{0.49\textwidth}
\begin{itemize}
\item Aquecimento extra de cabos e transformadores.
\item Sobrecarga do neutro por harmônicas triplas.
\item Perdas adicionais em motores.
\item Erros/interferências em equipamentos sensíveis.
\end{itemize}
\column{0.49\textwidth}
\begin{itemize}
\item Ressonância com bancos de capacitores.
\item Correntes elevadas em capacitores.
\item Atuação indevida de proteção.
\item Redução de vida útil de componentes.
\end{itemize}
\end{columns}
\begin{caixaProjeto}[title=Mitigação]
Reatores, filtros passivos/ativos, transformadores adequados, arquitetura de alimentação, seleção de conversores e projeto do banco de capacitores -- sempre após medição/estudo.
\end{caixaProjeto}
\end{frame}

\begin{frame}{Correção de fator de potência}
\small
\begin{columns}[T,onlytextwidth]
\column{0.51\textwidth}
\begin{align*}
Q_c=P\left(\tg\varphi_1-\tg\varphi_2\right)
\end{align*}
\begin{itemize}
\item Reduz corrente reativa circulante a montante do ponto de compensação.
\item Pode liberar capacidade e reduzir perdas.
\item Não corrige distorção harmônica.
\end{itemize}
\column{0.47\textwidth}
\begin{caixaEx}
$P=\SI{100}{kW}$, $fp_1=0{,}78$, $fp_2=0{,}95$:
\[
Q_c\approx\SI{47}{kvar}.
\]
Em ambiente com VFD, verificar ressonância e considerar banco dessintonizado/filtro conforme estudo.
\end{caixaEx}
\end{columns}
\end{frame}

\begin{frame}{Gestão de demanda}
\small
\centering
\begin{tikzpicture}
\begin{axis}[width=10.6cm,height=4.5cm,xlabel={hora},ylabel={kW},xmin=0,xmax=24,ymin=0,ymax=110,grid=both,legend pos=north west]
\addplot[corPrim,thick,smooth] coordinates {(0,25)(4,20)(8,65)(10,80)(12,72)(14,88)(17,100)(18,96)(20,62)(24,32)};\addlegendentry{sem gestão}
\addplot[corAccent,thick,dashed,smooth] coordinates {(0,25)(4,20)(8,62)(10,72)(12,69)(14,75)(17,78)(18,76)(20,60)(24,32)};\addlegendentry{com gestão}
\end{axis}
\end{tikzpicture}
\begin{caixaObs}
Deslocar cargas flexíveis, limitar recarga de VE e coordenar climatização pode reduzir pico sem reduzir energia útil do processo.
\end{caixaObs}
\end{frame}

\begin{frame}{Medição: o que registrar}
\scriptsize
\begin{tabularx}{\textwidth}{>{\bfseries}p{4.0cm}X}
\toprule
Grandeza & Uso \\
\midrule
$V$, $I$, $P$, $Q$, $S$, fp & Estado operacional e balanceamento. \\
Demanda e energia & Custos, contratos e eficiência. \\
THD/espectro & Diagnóstico de cargas não lineares. \\
Afundamentos e elevações & Falhas de processo e qualidade de suprimento. \\
Interrupções & Continuidade e confiabilidade. \\
Temperatura e termografia & Conexões, sobrecarga e ventilação. \\
Eventos de proteção & Causa-raiz e seletividade. \\
\bottomrule
\end{tabularx}
\end{frame}

%=====================================================================


% Cargas modernas prediais e documentação
\section{Veículos elétricos e cargas modernas}

\begin{frame}{Recarga de VE: por que não é uma tomada comum}
\small
\begin{columns}[T,onlytextwidth]
\column{0.50\textwidth}
\begin{itemize}
\item Potência elevada durante longos períodos.
\item Carga eletrônica com requisitos próprios de proteção.
\item Possibilidade de corrente residual CC conforme arquitetura.
\item Expansão simultânea de vários carregadores.
\item Aumento de carga pode alterar o padrão de entrada.
\item Gestão de demanda pode se tornar parte do projeto.
\end{itemize}
\column{0.48\textwidth}
\begin{caixaNorma}[title=Referências]
\begin{itemize}
\item ABNT NBR 17019:2022 — instalação fixa para alimentação/recepção de energia de VE;
\item NT.00042.EQTL Rev.04 — conexão de estações de recarga na área da Equatorial;
\item NT.00001/NT.00004 — entrada individual ou empreendimento com múltiplas UCs, conforme o caso.
\end{itemize}
\end{caixaNorma}
\end{columns}
\end{frame}

\begin{frame}{Arquitetura correta de um ponto de recarga}
\small
\centering
\begin{tikzpicture}[scale=0.78,every node/.style={font=\scriptsize}]
\tikzset{b/.style={draw=corPrim,fill=corDef!7,rounded corners,minimum width=2.25cm,minimum height=0.68cm,align=center}}
\node[b] (q) at (0,2.1) {QD};
\node[b] (p) at (3.0,2.1) {Proteção dedicada\\sobrecorrente + DR};
\node[b] (e) at (6.0,2.1) {SAVE / EVSE};
\node[b] (v) at (9.0,2.1) {Veículo};
\node[b] (pe) at (4.5,0.35) {Barramento PE / BEP\\equipotencialização};
\draw[-{Stealth},thick,corPrim] (q)--(p)--(e)--(v);
\draw[corNorma,thick] (q.south)--(pe.west);
\draw[corNorma,thick] (p.south)--(pe);
\draw[corNorma,thick] (e.south)--(pe.east);
\draw[corNorma,thick] (pe.south)--(4.5,-0.45) node[ground]{};
\end{tikzpicture}
\vspace{2mm}
\begin{caixaAt}
O PE é contínuo e integrado ao sistema de equipotencialização. Não se desenha um eletrodo de terra independente para cada equipamento. O tipo de DR deve ser compatível com a corrente residual possível, a detecção CC incorporada ao SAVE e as instruções normativas/do fabricante.
\end{caixaAt}
\end{frame}

\begin{frame}{Exemplo: SAVE monofásico de 7,4 kW}
\small
\begin{columns}[T,onlytextwidth]
\column{0.49\textwidth}
\begin{align*}
P&=\SI{7.4}{kW},\quad V=\SI{220}{V}\\
I_b&\approx\frac{7400}{220}=\SI{33.6}{A}
\end{align*}
\begin{itemize}
\item circuito dedicado;
\item cabo por ampacidade + queda + método;
\item proteção contra sobrecorrente;
\item proteção diferencial compatível;
\item PE e DPS conforme o sistema.
\end{itemize}
\column{0.49\textwidth}
\begin{caixaProjeto}[title=Impacto no projeto residencial]
No exemplo-base deste material:
\[
33{,}96+7{,}40=\SI{41.36}{kW}.
\]
Logo, a inclusão do SAVE não é apenas ``mais um circuito'': no exemplo, a faixa da Tabela 1 da NT.00001 passa de 24,1--38 kW (63 A) para 38,1--48 kW (80 A). QD, alimentador, balanceamento e solicitação de aumento de carga devem ser revistos.
\end{caixaProjeto}
\end{columns}
\end{frame}

\begin{frame}{Gestão dinâmica de recarga}
\small
\centering
\begin{tikzpicture}[scale=0.80,every node/.style={font=\scriptsize}]
\tikzset{b/.style={draw=corPrim,fill=corDef!7,rounded corners,minimum width=2.25cm,minimum height=0.65cm,align=center}}
\node[b] (med) at (0,2.2) {Medição da demanda\\da instalação};
\node[b] (ctrl) at (3.2,2.2) {Controlador\\de carga};
\node[b] (ev1) at (6.4,3.2) {SAVE 1};
\node[b] (ev2) at (6.4,2.2) {SAVE 2};
\node[b] (ev3) at (6.4,1.2) {SAVE 3};
\node[b] (pred) at (0,0.5) {Demais cargas\\da edificação};
\draw[-{Stealth},thick,corPrim] (med)--(ctrl);
\foreach \x in {ev1,ev2,ev3}{\draw[-{Stealth},thick,corAccent] (ctrl)--(\x);}
\draw[-{Stealth},thick,corPrim] (med)--(pred);
\node[align=center] at (3.2,0.45) {limite disponível\\+ prioridades};
\end{tikzpicture}
\vspace{1mm}
\begin{caixaObs}
Gestão de carga reduz simultaneidade dos carregadores, mas não elimina as verificações de capacidade térmica, proteção, qualidade de energia e requisitos da distribuidora.
\end{caixaObs}
\end{frame}

\begin{frame}{Geração distribuída: interface com a instalação}
\small
\begin{columns}[T,onlytextwidth]
\column{0.49\textwidth}
\begin{itemize}
\item Corrente pode fluir em sentido diferente do projeto tradicional.
\item Proteções e seccionamento existem nos lados CC e CA.
\item DPS e equipotencialização precisam ser compatíveis com o SPDA.
\item Verificar capacidade de barramentos e condutores no ponto de conexão.
\item O estudo interno deve fechar com a solicitação de conexão à distribuidora.
\end{itemize}
\column{0.49\textwidth}
\begin{caixaNorma}[title=Equatorial]
A conexão de micro e minigeração distribuída é tratada pela \textbf{NT.00020.EQTL Rev.06}. Para unidades do Grupo B, usar também os formulários e anexos atuais disponibilizados pela distribuidora.
\end{caixaNorma}
\begin{caixaAt}
Geração distribuída não autoriza reduzir proteção, seção de condutor ou capacidade de barramento apenas porque parte da energia é produzida localmente.
\end{caixaAt}
\end{columns}
\end{frame}

\begin{frame}{UPS, geradores e transferência}
\small
\begin{columns}[T,onlytextwidth]
\column{0.49\textwidth}
\begin{itemize}
\item Definir cargas essenciais e não essenciais.
\item Analisar curto-circuito em modo rede e modo fonte alternativa.
\item A proteção que atua bem na rede pode não atuar igual com fonte limitada eletronicamente.
\item Verificar neutro, aterramento e chaveamento conforme a topologia.
\end{itemize}
\column{0.49\textwidth}
\centering
\begin{tikzpicture}[scale=0.78,every node/.style={font=\scriptsize}]
\tikzset{b/.style={draw=corPrim,fill=corDef!7,rounded corners,minimum width=2.1cm,minimum height=0.62cm,align=center}}
\node[b] (r) at (0,2.6) {Rede};
\node[b] (g) at (0,0.8) {Gerador};
\node[b] (ats) at (3.2,1.7) {ATS};
\node[b] (ups) at (6.0,1.7) {UPS};
\node[b] (crit) at (8.8,1.7) {Cargas críticas};
\draw[-{Stealth},thick,corPrim] (r)--(ats);\draw[-{Stealth},thick,corAccent] (g)--(ats);\draw[-{Stealth},thick,corPrim] (ats)--(ups)--(crit);
\end{tikzpicture}
\end{columns}
\end{frame}

%=====================================================================

\section{CAD, BIM, simulação e documentação}

\begin{frame}{Ferramentas: cada uma responde a uma pergunta}
\scriptsize
\begin{tabularx}{\textwidth}{>{\bfseries}p{2.6cm}X}
\toprule
Ferramenta & Uso \\
\midrule
CAD elétrico & Plantas, diagramas, detalhes, listas e documentação 2D. \\
BIM & Compatibilização espacial, parâmetros, quantitativos e coordenação multidisciplinar. \\
QElectroTech & Diagramas didáticos/funcionais de comando e potência. \\
ETAP / EasyPower / similares & Fluxo de carga, curto, seletividade, partida, arco e harmônicas conforme módulos. \\
Planilha / Python & Automação de cálculos, auditoria, geração de tabelas e checagens repetitivas. \\
Analisador de energia & Validar modelo com medições reais em retrofit/operação. \\
\bottomrule
\end{tabularx}
\end{frame}

\begin{frame}{Fluxo de simulação elétrica}
\small
\centering
\begin{tikzpicture}[node distance=0.30cm,every node/.style={font=\scriptsize}]
\tikzset{b/.style={draw=corPrim,fill=corDef!7,rounded corners,minimum width=2.35cm,minimum height=0.65cm,align=center}}
\node[b] (a) {Topologia e fontes};
\node[b,right=of a] (b) {Cabos / trafos};
\node[b,right=of b] (c) {Cargas / modos};
\node[b,below=of c] (d) {Fluxo de carga};
\node[b,left=of d] (e) {Curto-circuito};
\node[b,left=of e] (f) {Seletividade / arco};
\node[b,below=of f] (g) {Partida / harmônicas};
\node[b,right=of g] (h) {Relatórios};
\node[b,right=of h] (i) {Revisão do projeto};
\foreach \u/\v in {a/b,b/c,c/d,d/e,e/f,f/g,g/h,h/i}{\draw[-{Stealth},thick,corPrim] (\u)--(\v);}
\draw[-{Stealth},thick,corAccent] (i.east) to[out=0,in=-80] (a.south east);
\end{tikzpicture}
\vspace{2mm}
\begin{caixaAt}
Software calcula o modelo fornecido. Dado errado + modelo errado = resultado preciso e inútil.
\end{caixaAt}
\end{frame}

\begin{frame}{BIM: parâmetros elétricos úteis}
\small
\begin{columns}[T,onlytextwidth]
\column{0.49\textwidth}
\begin{itemize}
\item Circuito e quadro de origem.
\item Potência e tensão.
\item Número de fases.
\item Fator de potência e demanda.
\item Reserva e status do circuito.
\end{itemize}
\column{0.49\textwidth}
\begin{itemize}
\item Comprimento e rota.
\item Cabo/eletroduto/leito.
\item Dispositivo de proteção.
\item Classificação de carga crítica.
\item Código para operação/manutenção.
\end{itemize}
\end{columns}
\begin{caixaObs}
O valor do BIM aumenta quando os parâmetros alimentam cálculo, documentação e operação -- não apenas um modelo 3D bonito.
\end{caixaObs}
\end{frame}

\begin{frame}[fragile]{Automação de checagens com Python}
\scriptsize
\begin{lstlisting}[style=codPy]
from math import sqrt, acos, sin

def corrente_3f(P, Vll, fp=1.0, eta=1.0):
    return P/(sqrt(3)*Vll*fp*eta)

def queda_3f(I, L_m, R_ohm_km, X_ohm_km, fp, Vll):
    phi = acos(fp)
    dv = sqrt(3)*I*(L_m/1000)*(R_ohm_km*fp + X_ohm_km*sin(phi))
    return 100*dv/Vll

I = corrente_3f(42000, 380, 0.88)
dv = queda_3f(I, 60, 0.75, 0.08, 0.88, 380)
print(I, dv)
\end{lstlisting}
\begin{caixaObs}
O script deve ler dados rastreáveis e gerar alerta, não escolher cabo ``por mágica''. O engenheiro continua responsável pelas premissas e validação.
\end{caixaObs}
\end{frame}

\begin{frame}{Documentação para manutenção}
\small
\begin{columns}[T,onlytextwidth]
\column{0.49\textwidth}
\begin{itemize}
\item Unifilar atualizado.
\item Ajustes de relés/disjuntores.
\item Curto-circuito disponível por barra.
\item Identificação de fontes alternativas.
\item Procedimentos de manobra.
\end{itemize}
\column{0.49\textwidth}
\begin{itemize}
\item Etiquetas e riscos aplicáveis.
\item Histórico de intervenções.
\item Ensaios de DR, isolamento e aterramento.
\item Termografias e inspeções.
\item Lista de sobressalentes críticos.
\end{itemize}
\end{columns}
\end{frame}

%=====================================================================


% Projeto residencial completo + exemplos
%=====================================================================
% PROJETO RESIDENCIAL I — CASA TERREA COM DUAS SUITES
% Base eletrica: ABNT NBR 5410 + NT.00001.EQTL Rev.09
% Planta-base didatica com logica arquitetonica e mobiliario.
% Atualizacao: agosto de 2026
%=====================================================================

\ifdefined\ArqParedeExt\else\input{projetos/base/plantas_arquitetonicas.tex}\fi

\section{Projeto Residencial I — casa térrea com duas suítes}

% Símbolos elétricos locais
\providecommand{\PRLuz}[2]{\begin{scope}[shift={(#1,#2)}]\draw[corAccent,thick] (0,0) circle (0.14);\draw[corAccent,thick] (-0.10,-0.10)--(0.10,0.10);\draw[corAccent,thick] (-0.10,0.10)--(0.10,-0.10);\end{scope}}
\providecommand{\PRTUG}[2]{\begin{scope}[shift={(#1,#2)}]\draw[corPrim,thick] (0,0) circle (0.13);\draw[corPrim,thick] (-0.16,0)--(0.16,0);\end{scope}}
\providecommand{\PRTUE}[3]{\begin{scope}[shift={(#1,#2)}]\draw[corAt,thick] (0,0) circle (0.17);\node[font=\tiny\bfseries,corAt] at (0,0){#3};\end{scope}}
\providecommand{\PRQD}[3]{\begin{scope}[shift={(#1,#2)}]\draw[corPrim,fill=corDef!10,thick,rounded corners=1pt] (-0.38,-0.22) rectangle (0.38,0.22);\node[font=\tiny\bfseries] at (0,0){#3};\end{scope}}

%---------------------------------------------------------------------
% Planta 15,00 x 11,00 m. Frente em y=0.
% Ala íntima à esquerda; social ao centro; garagem/serviço à direita.
%---------------------------------------------------------------------
\newcommand{\CasaResidencialBase}{%
  % contorno
  \draw[arqExt] (0,0) rectangle (15,11);
  % ala íntima e hall
  \ArqParedeInt{5.2}{0}{5.2}{11}
  \ArqParedeInt{6.4}{0}{6.4}{11}
  \ArqParedeInt{0}{3.6}{5.2}{3.6}
  \ArqParedeInt{0}{7.2}{5.2}{7.2}
  % banheiros internos / social
  \ArqParedeInt{3.5}{0}{3.5}{2.0}
  \ArqParedeInt{3.5}{2.0}{5.2}{2.0}
  \ArqParedeInt{3.5}{4.8}{5.2}{4.8}
  \ArqParedeInt{3.5}{4.8}{3.5}{7.2}
  \ArqParedeInt{3.5}{8.4}{5.2}{8.4}
  \ArqParedeInt{3.5}{8.4}{3.5}{11}
  % social/cozinha
  \ArqParedeInt{6.4}{6.4}{11.2}{6.4}
  % garagem / serviço / varanda gourmet
  \ArqParedeInt{11.2}{0}{11.2}{11}
  \ArqParedeInt{11.2}{5.6}{15}{5.6}
  \ArqParedeInt{11.2}{8.2}{15}{8.2}

  % acessos: principal diretamente na sala; garagem para sala; serviço externo
  \ArqPortaHUp{7.25}{0}{1.10}
  \draw[white,line width=5pt] (11.2,2.2)--(11.2,3.2);\draw[black!60,thin] (11.2,2.2)--(10.2,2.2);
  \ArqPortaVLeft{15}{6.25}{0.90}
  \draw[white,line width=5pt] (11.2,9.1)--(11.2,10.7);\draw[arqGlass] (11.2,9.1)--(11.2,10.7);
  % hall interno aberto para a sala
  \draw[white,line width=5pt] (6.4,2.45)--(6.4,3.65);
  % portas dos quartos/banheiros para hall
  \ArqPortaVLeft{5.2}{2.45}{0.90}
  \ArqPortaVLeft{5.2}{3.85}{0.90}
  \ArqPortaVLeft{5.2}{7.45}{0.90}
  \ArqPortaVLeft{5.2}{0.55}{0.80}
  % banheiros das suítes: acesso exclusivo por dentro do quarto
  \ArqPortaVLeft{3.5}{5.05}{0.80}
  \ArqPortaVLeft{3.5}{8.65}{0.80}
  % cozinha -> serviço
  \ArqPortaVRight{11.2}{6.75}{0.85}
  % vão de garagem
  \draw[white,line width=5.5pt] (11.55,0)--(14.70,0);\draw[black!30,very thin] (11.55,0.10)--(14.70,0.10);

  % janelas
  \ArqJanV{0.45}{2.75}{0}
  \ArqJanV{4.15}{6.55}{0}
  \ArqJanV{7.75}{10.20}{0}
  \ArqJanH{8.0}{10.2}{0}
  \ArqJanH{7.3}{9.8}{11}
  \ArqJanH{0.7}{2.6}{11}
  \ArqJanH{12.0}{14.1}{11}

  % mobiliário - quarto 3
  \ArqCama{0.55}{0.45}{2.15}{2.65}
  \ArqArmario{2.85}{2.45}{0.55}{0.95}
  \node[arqLabel] at (1.75,3.25){QUARTO 3};
  % banho social
  \ArqChuveiro{3.65}{0.15}{1.25}{0.85}\ArqVaso{4.55}{1.35}\ArqPia{3.62}{1.20}{0.65}{0.45}
  \node[arqSub] at (4.30,1.82){BANHO SOCIAL};

  % suíte 2
  \ArqCama{0.55}{4.05}{2.20}{2.65}\ArqArmario{2.85}{3.95}{0.50}{1.05}
  \ArqChuveiro{3.65}{6.05}{1.25}{0.95}\ArqVaso{4.55}{5.35}\ArqPia{3.62}{5.05}{0.65}{0.45}
  \node[arqLabel] at (1.75,6.85){SUÍTE 2};\node[arqSub] at (4.35,6.90){B2};

  % suíte master
  \ArqCama{0.55}{7.75}{2.35}{2.70}\ArqArmario{2.95}{7.45}{0.45}{0.85}
  \ArqChuveiro{3.65}{10.00}{1.25}{0.85}\ArqVaso{4.55}{9.20}\ArqPia{3.62}{8.65}{0.65}{0.45}
  \node[arqLabel] at (1.85,10.65){SUÍTE MASTER};\node[arqSub] at (4.35,10.70){B1};
  \node[arqSub,rotate=90] at (5.80,5.50){HALL ÍNTIMO};

  % sala/jantar
  \ArqSofa{7.00}{0.85}{3.10}{1.05}\ArqBancada{10.50}{0.95}{0.35}{2.25}
  \ArqMesa{7.25}{4.20}{2.55}{1.10}
  \node[arqLabel] at (8.80,3.20){SALA / JANTAR};
  % cozinha
  \ArqBancada{6.75}{10.15}{3.50}{0.55}\ArqBancada{10.35}{7.05}{0.55}{3.65}
  \ArqFogao{9.25}{10.05}\ArqGeladeira{10.15}{6.65}\ArqBancada{7.40}{7.55}{2.30}{0.80}
  \node[arqLabel] at (8.60,9.45){COZINHA};
  % garagem
  \ArqCarro{12.00}{0.65}{2.15}{4.25}\node[arqLabel] at (13.10,5.20){GARAGEM};
  % serviço
  \ArqLavadora{12.00}{6.05}\ArqBancada{13.0}{6.0}{1.45}{0.55}\node[arqLabel] at (13.10,7.65){SERVIÇO};
  % gourmet / quintal coberto
  \ArqBancada{12.0}{9.85}{2.35}{0.55}\ArqMesa{12.20}{8.70}{1.55}{0.75}\node[arqLabel] at (13.10,10.75){VARANDA GOURMET};

  \CotaH{0}{15}{11.18}{15,00 m}\CotaV{0}{11}{15.18}{11,00 m}
}

\begin{frame}{Residencial I — planta arquitetônica funcional}
\centering
\begin{tikzpicture}[scale=0.49,every node/.style={font=\tiny}]
\CasaResidencialBase
\end{tikzpicture}
\vspace{0.5mm}
\scriptsize A entrada principal chega à sala. O hall é interno. A casa possui três banheiros completos: banho social, B2 e B1; estes dois últimos são acessados exclusivamente pelas suítes.
\end{frame}

\begin{frame}{Residencial I — iluminação e circuitos C1/C2}
\centering
\begin{tikzpicture}[scale=0.49,every node/.style={font=\tiny}]
\CasaResidencialBase
\PRQD{6.10}{3.10}{QD1}
\foreach \x/\y in {1.7/2.1,4.25/1.0,1.75/5.4,4.25/5.9,1.8/9.2,4.25/9.6,5.8/5.5,8.8/2.8,8.7/5.2,8.6/8.8,13.1/3.0,13.1/6.9,13.1/9.6}{\PRLuz{\x}{\y}}
% C1: sala, cozinha, servico e gourmet
\draw[corAccent!80!black,dashed,thick] (6.10,3.10)--(8.8,2.8)--(8.7,5.2)--(8.6,8.8)--(13.1,9.6);
\draw[corAccent!80!black,dashed,thick] (8.6,8.8)--(13.1,6.9);
% C2: ala intima, banheiros, hall e garagem
\draw[corAccent!55!black,dashed,thick] (6.10,3.10)--(5.8,5.5)--(1.75,5.4)--(4.25,5.9)--(1.8,9.2)--(4.25,9.6);
\draw[corAccent!55!black,dashed,thick] (6.10,3.10)--(1.7,2.1)--(4.25,1.0);
\draw[corAccent!55!black,dashed,thick] (6.10,3.10)--(13.1,3.0);
\node[font=\tiny\bfseries,fill=white,inner sep=1pt,corAccent!80!black] at (9.5,7.4){C1};
\node[font=\tiny\bfseries,fill=white,inner sep=1pt,corAccent!55!black] at (3.0,5.4){C2};
\end{tikzpicture}
\vspace{0.4mm}
\scriptsize Previsão calculada: C1 = 1,06 kVA e C2 = 1,40 kVA. Os 13 símbolos são pontos de luz; os 2,46 kVA resultam das áreas e dos acréscimos completos de 4 m$^2$, não da quantidade de símbolos.
\end{frame}

\begin{frame}{Residencial I — 43 TUG e oito cargas específicas}
\centering
\begin{tikzpicture}[scale=0.49,every node/.style={font=\tiny}]
\CasaResidencialBase
\PRQD{6.10}{3.10}{QD1}
% C3: sala e dormitorios — 17 pontos
\foreach \x/\y in {0.25/0.35,0.25/3.30,3.25/0.35,3.25/3.30,0.25/3.85,0.25/6.95,3.25/3.85,3.25/6.95,0.25/7.45,0.25/10.75,3.25/7.45,3.25/10.75,6.65/0.35,10.95/0.35,10.95/3.60,10.95/6.15,6.65/6.15}{\PRTUG{\x}{\y}}
% C4: hall e garagem — 9 pontos
\foreach \x/\y in {6.25/0.70,6.25/2.10,6.25/4.20,6.25/7.80,6.25/10.20,11.35/0.80,11.35/4.80,14.65/1.50,14.65/4.80}{\PRTUG{\x}{\y}}
% C5 cozinha; C6 servico; C7 gourmet; C8 banheiros
\foreach \x/\y in {6.70/10.75,8.00/10.75,9.30/10.75,10.95/7.00,10.95/8.50,10.95/10.00}{\PRTUG{\x}{\y}}
\foreach \x/\y in {11.35/5.85,14.65/5.85,11.35/7.95,14.65/7.95}{\PRTUG{\x}{\y}}
\foreach \x/\y in {11.35/8.45,14.65/8.45,11.35/10.75,14.65/10.75}{\PRTUG{\x}{\y}}
\foreach \x/\y in {5.00/1.80,4.95/5.00,4.95/8.65}{\PRTUG{\x}{\y}}
% TUE C9--C16
\PRTUE{4.35}{0.55}{C9}\PRTUE{4.35}{6.55}{C10}\PRTUE{4.35}{10.45}{C11}
\PRTUE{9.60}{10.45}{C12}\PRTUE{12.35}{6.40}{C13}
\PRTUE{0.45}{2.55}{C14}\PRTUE{0.45}{6.25}{C15}\PRTUE{0.45}{9.95}{C16}
\end{tikzpicture}
\vspace{0.3mm}
\tiny C3 sala/quartos (17) · C4 hall/garagem (9) · C5 cozinha (6) · C6 serviço (4) · C7 gourmet (4) · C8 banheiros (3) · C9--C11 chuveiros · C12 forno · C13 lavadora · C14--C16 ar-condicionado.
\end{frame}

\begin{frame}[shrink=7]{Residencial I — memória dos mínimos por ambiente}
\scriptsize
\begin{tabularx}{\textwidth}{>{\bfseries}p{2.55cm}c c c c >{\centering\arraybackslash}X}
\toprule
Ambiente & Área (m$^2$) & Perím. (m) & Luz (VA) & TUG & TUG (VA)\\
\midrule
Quarto 3 & 15,3 & 17,6 & 220 & 4 & 400\\
Suíte 2 & 14,6 & 17,6 & 220 & 4 & 400\\
Suíte master & 15,3 & 18,0 & 220 & 4 & 400\\
Banheiros (3) & 3,4/4,1/4,4 & critério próprio & 300 & 3 & 1.800\\
Hall íntimo & 13,2 & 24,4 & 160 & 5 & 500\\
Sala/jantar & 30,7 & 22,4 & 460 & 5 & 500\\
Cozinha & 22,1 & 18,8 & 340 & 6 & 2.100\\
Garagem & 21,3 & 18,8 & 280 & 4 & 400\\
Serviço & 9,9 & 12,8 & 100 & 4 & 1.900\\
Varanda gourmet$^*$ & 10,6 & 13,2 & 160 & 4 & 1.900\\
\midrule
\textbf{Total} & -- & -- & \textbf{2.460} & \textbf{43} & \textbf{10.300}\\
\bottomrule
\end{tabularx}
\begin{caixaObs}
Áreas e perímetros seguem a geometria interna desenhada; os dormitórios têm forma em L. $^*$A varanda gourmet foi tratada como ambiente análogo a cozinha por possuir bancada de preparo.
\end{caixaObs}
\end{frame}

\begin{frame}[shrink=6]{Residencial I — circuitos C1 a C8}
\scriptsize
\begin{tabularx}{\textwidth}{c p{3.25cm} c c c c c X}
\toprule
C & Uso & $S$ & $I_b$ & Cu & DJ & Fase & Solução adotada\\
\midrule
C1 & Luz social/cozinha/serviço/gourmet & 1,06 kVA & 4,8 A & 1,5 mm$^2$ & 10 A & A & RCBO 30 mA\\
C2 & Luz íntima/banhos/hall/garagem & 1,40 kVA & 6,4 A & 1,5 mm$^2$ & 10 A & B & RCBO 30 mA\\
C3 & TUG sala e dormitórios & 1,70 kVA & 7,7 A & 2,5 mm$^2$ & 16 A & A & RCBO 30 mA\\
C4 & TUG hall e garagem & 0,90 kVA & 4,1 A & 2,5 mm$^2$ & 16 A & B & RCBO 30 mA\\
C5 & TUG cozinha & 2,10 kVA & 9,5 A & 2,5 mm$^2$ & 16 A & C & RCBO 30 mA\\
C6 & TUG serviço & 1,90 kVA & 8,6 A & 2,5 mm$^2$ & 16 A & C & RCBO 30 mA\\
C7 & TUG gourmet & 1,90 kVA & 8,6 A & 2,5 mm$^2$ & 16 A & A & RCBO 30 mA\\
C8 & TUG banheiros & 1,80 kVA & 8,2 A & 2,5 mm$^2$ & 16 A & C & RCBO 30 mA\\
\bottomrule
\end{tabularx}
\begin{caixaObs}
RCBO individuais evitam neutros compartilhados a jusante e limitam a área afetada. Tipo A é a premissa inicial; confirmar forma de onda residual, corrente de fuga e fabricante.
\end{caixaObs}
\end{frame}

\begin{frame}[shrink=6]{Residencial I — circuitos C9 a C16}
\scriptsize
\begin{tabularx}{\textwidth}{c p{3.15cm} c c c c c X}
\toprule
C & Equipamento & $P$ & $I_b$ & Cu & DJ & Fase & Proteção/observação\\
\midrule
C9 & Chuveiro banho social & 4,50 kW & 20,5 A & 4 mm$^2$ & 25 A & A & dedicado; RCBO 30 mA\\
C10 & Chuveiro suíte 2 & 4,50 kW & 20,5 A & 4 mm$^2$ & 25 A & B & dedicado; RCBO 30 mA\\
C11 & Chuveiro suíte master & 4,50 kW & 20,5 A & 4 mm$^2$ & 25 A & C & dedicado; RCBO 30 mA\\
C12 & Forno elétrico & 3,50 kW & 15,9 A & 4 mm$^2$ & 20 A & B & dedicado; RCBO 30 mA\\
C13 & Lavadora & 1,20 kW & 5,5 A & 2,5 mm$^2$ & 16 A & A & dedicado; RCBO 30 mA\\
C14--C16 & Ar-condicionado quartos & 3$\times$1,00 kW & 4,5 A cada & 2,5 mm$^2$ & 16 A & C/B/A & dedicados; fabricante\\
\bottomrule
\end{tabularx}
\begin{caixaAt}
As seções e proteções são seleções iniciais. Antes do executivo: método de instalação, agrupamento, temperatura, $I_z$, queda, curto máximo/mínimo, capacidade de interrupção e manual de cada equipamento.
\end{caixaAt}
\end{frame}

\begin{frame}[shrink=4]{Residencial I — entrada e balanceamento fechados}
\small
\begin{columns}[T,onlytextwidth]
\column{0.49\textwidth}
\begin{caixaProjeto}[title=Balanceamento por valores instalados]
\begin{align*}
S_A&=\SI{11.36}{kVA}\quad (\SI{51.6}{A})\\
S_B&=\SI{11.30}{kVA}\quad (\SI{51.4}{A})\\
S_C&=\SI{11.30}{kVA}\quad (\SI{51.4}{A})
\end{align*}
\scriptsize A: C1+C3+C7+C9+C13+C16\\
B: C2+C4+C10+C12+C15\\
C: C5+C6+C8+C11+C14
\end{caixaProjeto}
\column{0.49\textwidth}
\begin{caixaNorma}[title=Equatorial Maranhão]
Hipótese declarada do estudo: \SI{33.96}{kW} equivalentes. Na Tabela 1 da NT.00001.EQTL Rev.09:
\begin{itemize}
\item faixa 24,1--38 kW;
\item atendimento 380/220 V, 3F+N;
\item disjuntor de entrada 63 A;
\item condutor Cu do cliente no padrão: mínimo 10 mm$^2$.
\end{itemize}
\end{caixaNorma}
\end{columns}
\begin{caixaAt}
Para ilustrar o enquadramento, os VA mínimos de luz/TUG foram tomados numericamente como W. Na ligação real, declarar as potências ativas nominais efetivamente instaladas conforme a definição da NT.00001; demanda não substitui carga instalada nessa tabela. O alimentador interno é recalculado pela NBR 5410.
\end{caixaAt}
\end{frame}

\begin{frame}{Residencial I — unifilar coerente com 16 circuitos}
\centering
\begin{tikzpicture}[scale=0.76,every node/.style={font=\tiny}]
\tikzset{b/.style={draw=corPrim,fill=corDef!7,rounded corners,minimum width=2.05cm,minimum height=0.58cm,align=center}}
\node[b](r)at(0,4.6){Rede EQTL\\380/220 V};
\node[b](m)at(0,3.70){Medição\\padrão de entrada};
\node[b](dg)at(0,2.80){DG 3P 63 A\\+ DPS coordenado};
\node[b](qd)at(0,1.85){QD1\\barras N e PE separadas};
\node[b](g1)at(-4.2,0.45){C1--C8\\luz/TUG + RCBO};
\node[b](g2)at(-1.4,0.45){C9--C11\\chuveiros + RCBO};
\node[b](g3)at(1.4,0.45){C12--C13\\forno/lavadora};
\node[b](g4)at(4.2,0.45){C14--C16\\ar-condicionado};
\node[b,draw=corNorma,fill=corNorma!7](bep)at(5.0,2.65){BEP\\PE + equipot.};
\draw[-{Stealth},thick,corPrim](r)--(m)--(dg)--(qd);
\foreach \x in {g1,g2,g3,g4}{\draw[-{Stealth},thick](qd)--(\x);}
\draw[corNorma,thick](qd)--(bep);
\draw[corNorma,thick](bep.south)--(5.0,1.35)node[ground]{};
\node[font=\tiny,corNorma,anchor=west] at (5.15,1.35){eletrodo(s)};
\end{tikzpicture}
\vspace{1mm}
\scriptsize O PE não é um ``terra isolado'' por circuito; todos os circuitos retornam ao barramento PE, ligado ao BEP e ao sistema comum de aterramento/equipotencialização.
\end{frame}

\section{Exemplos resolvidos}

\begin{frame}{Exemplo 1 — circuito de TUG em área seca}
\small
\begin{columns}[T,onlytextwidth]
\column{0.49\textwidth}
\begin{caixaEx}
Circuito em 220 V, $P=\SI{1100}{VA}$, comprimento de estudo $L=\SI{30}{m}$:
\[
I_b=\frac{1100}{220}=\SI{5.00}{A}.
\]
Adota-se inicialmente Cu 2,5 mm$^2$ e disjuntor 16 A.
\end{caixaEx}
\column{0.49\textwidth}
\textbf{A seleção só fecha após verificar:}
\begin{enumerate}
\item método de instalação;
\item temperatura e agrupamento;
\item $I_b\le I_n\le I_z$;
\item queda de tensão;
\item curto e desligamento;
\item necessidade de proteção diferencial.
\end{enumerate}
\end{columns}
\end{frame}

\begin{frame}{Exemplo 2 — circuito de chuveiro}
\small
\begin{columns}[T,onlytextwidth]
\column{0.49\textwidth}
\begin{caixaEx}
$P=\SI{5.5}{kW}$, $V=\SI{220}{V}$:
\[
I_b=\frac{5500}{220}=\SI{25}{A}.
\]
Para este exemplo de 5,5 kW, adota-se inicialmente circuito dedicado, Cu 6 mm$^2$ e DJ 32 A, após as verificações de capacidade de condução e queda. O projeto residencial adiante usa chuveiros de 4,5 kW e outra seleção.
\end{caixaEx}
\column{0.49\textwidth}
\begin{caixaNorma}[title=Proteção adicional]
O circuito que serve ponto de utilização em local contendo chuveiro deve receber proteção diferencial-residual de alta sensibilidade conforme os requisitos aplicáveis.
\end{caixaNorma}
\begin{caixaAt}
Não escolher seção apenas por uma tabela de potência de chuveiro. Comprimento, método, temperatura e agrupamento mudam o resultado.
\end{caixaAt}
\end{columns}
\end{frame}

\begin{frame}{Exemplo 3 — balanceamento antes do padrão de entrada}
\small
\begin{columns}[T,onlytextwidth]
\column{0.52\textwidth}
No projeto residencial desenvolvido neste material:
\[
S_A=\SI{11.36}{kVA},\qquad S_B=S_C=\SI{11.30}{kVA}.
\]
Em 220 V fase-neutro:
\[
I_A\approx\SI{51.6}{A},\qquad I_B=I_C\approx\SI{51.4}{A}.
\]
\column{0.46\textwidth}
\begin{caixaProjeto}[title=Distribuição]
\begin{itemize}
\item A: C1+C3+C7+C9+C13+C16;
\item B: C2+C4+C10+C12+C15;
\item C: C5+C6+C8+C11+C14.
\end{itemize}
A soma fecha com os 16 circuitos e mantém um chuveiro em cada fase. Este é o balanceamento dos valores instalados; cenários de demanda também devem ser verificados.
\end{caixaProjeto}
\end{columns}
\end{frame}

\begin{frame}{Exemplo 4 — enquadramento Equatorial Maranhão}
\small
\begin{columns}[T,onlytextwidth]
\column{0.50\textwidth}
Carga instalada do projeto:
\[
P_{inst,eq}=\SI{33.96}{kW}.
\]
A NT.00001.EQTL Rev.09 enquadra a faixa de 24,1 a 38 kW, na Tabela 1 de 220/380 V, em ligação trifásica com disjuntor termomagnético de entrada de 63 A.
\column{0.48\textwidth}
\begin{caixaNorma}[title=No Maranhão]
A ligação trifásica urbana é 380/220 V, 3F+N. A mesma tabela indica 10 mm$^2$ de cobre como mínimo do condutor do cliente no padrão correspondente.
\end{caixaNorma}
\begin{caixaAt}
O alimentador interno é recalculado pela NBR 5410; não se copia automaticamente o mínimo do padrão da concessionária.
\end{caixaAt}
\end{columns}
\end{frame}

\begin{frame}{Exemplo 5 — recarga de VE e aumento de carga}
\small
\begin{columns}[T,onlytextwidth]
\column{0.49\textwidth}
Se ao projeto de $\SI{33.96}{kW}$ for acrescentado um SAVE de $\SI{7.4}{kW}$:
\[
P_{inst,novo}=33{,}96+7{,}40=\SI{41.36}{kW}.
\]
O exemplo passa à faixa de 38,1 a 48 kW, associada a 80 A na Tabela 1; o enquadramento da entrada precisa ser revisto.
\column{0.49\textwidth}
\begin{caixaProjeto}[title=Consequência]
Além do novo circuito e da capacidade do QD/alimentador, a conexão da estação de recarga deve observar a NT.00042.EQTL Rev.04 e eventual necessidade de aumento de carga perante a distribuidora.
\end{caixaProjeto}
\end{columns}
\end{frame}

%=====================================================================

\section{Exemplos resolvidos}

\begin{frame}{Exemplo 2 -- alimentador trifásico}
\small
\begin{columns}[T,onlytextwidth]
\column{0.50\textwidth}
\begin{caixaEx}
$P=\SI{42}{kW}$, $V=\SI{380}{V}$, $fp=0{,}88$, $L=\SI{60}{m}$.
\[
I_b=\frac{42000}{\sqrt3\cdot380\cdot0{,}88}\approx\SI{72.5}{A}.
\]
\end{caixaEx}
\column{0.48\textwidth}
\begin{caixaProjeto}[title=Sequência]
\begin{enumerate}
\item selecionar seção por ampacidade corrigida;
\item selecionar proteção;
\item verificar $\Delta V$;
\item verificar curto máximo/mínimo;
\item verificar PE e neutro;
\item revisar seletividade a montante/jusante.
\end{enumerate}
\end{caixaProjeto}
\end{columns}
\end{frame}
\begin{frame}{Exemplo 3 -- banco de capacitores}
\small
\begin{columns}[T,onlytextwidth]
\column{0.50\textwidth}
\begin{caixaEx}
$P=\SI{60}{kW}$, corrigir $fp$ de $0{,}80$ para $0{,}95$.
\begin{align*}
\varphi_1&=36{,}87^\circ,\quad \varphi_2=18{,}19^\circ\\
Q_c&=60(\tg\varphi_1-\tg\varphi_2)\\
&\approx\SI{25.3}{kvar}.
\end{align*}
\end{caixaEx}
\column{0.48\textwidth}
\begin{caixaAt}
Antes de instalar $\SI{25}{kvar}$ ``puro'', avaliar:
\begin{itemize}
\item perfil de carga;
\item risco de sobrecompensação;
\item harmônicas;
\item ressonância;
\item número de estágios;
\item local de compensação.
\end{itemize}
\end{caixaAt}
\end{columns}
\end{frame}
\begin{frame}{Exemplo 4 -- motor de 22 kW em rede fraca}
\small
\begin{columns}[T,onlytextwidth]
\column{0.50\textwidth}
\begin{caixaEx}
$P=\SI{22}{kW}$, $V=\SI{380}{V}$, $fp=0{,}86$, $\eta=0{,}93$.
\[
I_N\approx\frac{22000}{\sqrt3\cdot380\cdot0{,}86\cdot0{,}93}
\approx\SI{41.8}{A}.
\]
\end{caixaEx}
\column{0.48\textwidth}
\begin{caixaProjeto}[title=Decisão]
Se a rede não tolera inrush elevado:
\begin{itemize}
\item calcular queda durante partida;
\item verificar torque requerido;
\item comparar soft-starter e VFD;
\item coordenar proteção;
\item verificar harmônicas/EMC.
\end{itemize}
\end{caixaProjeto}
\end{columns}
\end{frame}
\begin{frame}{Exemplo 5 -- curto no secundário do transformador}
\small
\begin{columns}[T,onlytextwidth]
\column{0.51\textwidth}
\begin{caixaEx}
Trafo $\SI{500}{kVA}$, $\SI{380}{V}$, $Z=5{,}5\%$.
\[
I_N=\frac{500000}{\sqrt3\cdot380}\approx\SI{760}{A}.
\]
\[
I_{cc}\approx\frac{760}{0{,}055}\approx\SI{13.8}{kA}.
\]
\end{caixaEx}
\column{0.47\textwidth}
\begin{caixaObs}
No QGBT real, incluir impedância da rede de MT, trafo, conexões e barramentos. Nos QDs a jusante, incluir cabos/alimentadores, reduzindo $I_{cc}$.
\end{caixaObs}
\end{columns}
\end{frame}
\begin{frame}{Exemplo 6 -- balanceamento de fases}
\small
\begin{columns}[T,onlytextwidth]
\column{0.50\textwidth}
\begin{caixaEx}
Correntes medidas:
\[
I_A=\SI{52}{A},\quad I_B=\SI{38}{A},\quad I_C=\SI{44}{A}.
\]
\[
I_{med}=\SI{44.7}{A}.
\]
A fase A está $\approx\SI{7.3}{A}$ acima da média.
\end{caixaEx}
\column{0.48\textwidth}
\begin{caixaProjeto}[title=Ação]
Mover circuitos monofásicos de A para B/C, mas respeitando:
\begin{itemize}
\item simultaneidade real;
\item cargas críticas;
\item sequência de fases;
\item DRs e neutros independentes;
\item harmônicas no neutro.
\end{itemize}
\end{caixaProjeto}
\end{columns}
\end{frame}



% Exercícios
\section{Exercícios N1--N4}

\begin{frame}{N1 — leitura da planta e requisitos mínimos}
\small
\begin{enumerate}
\item Para uma sala de 18 m$^2$ e perímetro de 17 m, determine a previsão mínima de carga de iluminação e o número mínimo de TUG.
\item Explique por que a potência de previsão de iluminação não determina diretamente a quantidade de luminárias.
\item Em uma planta residencial, indique os locais adequados para QD, interruptores e tomadas considerando portas, circulação e mobiliário.
\item Diferencie TUG, TUE, circuito dedicado, PE, DR/DDR e DPS.
\end{enumerate}
\end{frame}

\begin{frame}{N2 — dimensionamento de circuitos}
\small
\begin{enumerate}
\item Dimensione um circuito de TUG em 220 V, 1.900 VA e 24 m, verificando corrente, proteção, seção inicial e queda de tensão.
\item Analise um chuveiro de 5,5 kW em 220 V a 28 m do QD; verifique a seção por ampacidade e queda de tensão.
\item Proponha a divisão de circuitos para uma residência com sala, cozinha, área de serviço, dois quartos, banheiro e escritório.
\item Explique quais desses circuitos exigem proteção diferencial-residual de alta sensibilidade e por quê.
\end{enumerate}
\end{frame}

\begin{frame}{N3 — quadro, balanceamento e entrada Equatorial}
\small
\begin{enumerate}
\item Distribua 11 circuitos monofásicos de potências distintas entre A, B e C em uma rede 380/220 V, minimizando o desbalanceamento da carga instalada.
\item Para uma unidade residencial no Maranhão com 21,5 kW instalados, determine o enquadramento de fornecimento usando a NT.00001.EQTL Rev.09 e identifique quais dados ainda precisam ser verificados no padrão de entrada.
\item Mostre por que o condutor mínimo indicado pela concessionária para o padrão não pode ser copiado automaticamente para o alimentador interno.
\item Desenhe um unifilar com DG, DPS, barramentos N/PE, circuitos com DDR/RCBO e identificação coerente com o quadro de cargas.
\end{enumerate}
\end{frame}

\begin{frame}{N4 — desafio de projeto residencial completo}
\small
Receba uma planta arquitetônica residencial cotada e, sem alterar sua geometria, desenvolva o projeto elétrico completo:
\begin{enumerate}
\item calcule os pontos mínimos de iluminação e TUG e justifique pontos adicionais;
\item posicione TUE e cargas específicas conforme uso real;
\item divida e identifique os circuitos na planta;
\item dimensione condutores, eletrodutos e proteções;
\item verifique queda de tensão, proteção diferencial, DPS, PE e equipotencialização;
\item balanceie as fases e enquadre a entrada conforme Equatorial Maranhão;
\item gere quadro de cargas, unifilar, memorial de cálculo e lista de materiais;
\item confira se planta, quadro e unifilar apresentam exatamente os mesmos circuitos.
\end{enumerate}
\end{frame}

%=====================================================================


\section{Exercícios N1--N4}
\begin{frame}{N1 — fundamentos industriais}
\small
\begin{enumerate}
\item Diferencie QGBT, QD, CCM e painel de automação.
\item Explique as funções de contator, relé de sobrecarga, disjuntor-motor, soft-starter e VFD.
\item Calcule a corrente nominal de um motor trifásico de \SI{15}{kW}, \SI{380}{V}, $fp=0{,}86$ e $\eta=0{,}90$.
\item Descreva a função da equipotencialização em uma planta industrial.
\end{enumerate}
\end{frame}

\begin{frame}{N2 — dimensionamento industrial}
\small
\begin{enumerate}
\item Dimensione preliminarmente o alimentador de um motor de \SI{30}{kW} a \SI{55}{m} do CCM.
\item Compare partida direta, estrela-triângulo, soft-starter e VFD para uma carga de alto torque de partida.
\item Verifique queda de tensão de um alimentador trifásico e proponha duas correções possíveis.
\item Defina critérios para seleção da capacidade de interrupção de um disjuntor industrial.
\end{enumerate}
\end{frame}

\begin{frame}{N3 — análise industrial}
\small
\begin{enumerate}
\item Estime a corrente de curto-circuito no secundário de um transformador a partir de sua potência e impedância percentual.
\item Discuta como obter seletividade entre proteção do QGBT e do CCM.
\item Avalie os efeitos de VFDs sobre harmônicas, fator de potência e corrente no neutro.
\item Proponha ações para reduzir risco de arco elétrico durante manutenção.
\end{enumerate}
\end{frame}

\begin{frame}{N4 — desafio de projeto industrial}
\small
Uma planta possui transformador próprio, QGBT, dois CCMs, motores de \SI{45}{kW}, \SI{30}{kW} e \SI{18.5}{kW}, cargas de soldagem e painéis de automação.
\begin{enumerate}
\item Estruture a arquitetura elétrica e o diagrama unifilar.
\item Defina critérios de demanda, curto-circuito e seletividade.
\item Proponha alimentação, proteção e acionamento dos motores.
\item Inclua aterramento, qualidade de energia, NR-10 e plano de manutenção.
\end{enumerate}
\end{frame}



\section{Referências e rastreabilidade}

\begin{frame}[shrink=4]{Fontes oficiais verificadas nesta revisão}
\small
\begin{itemize}
\item \textbf{ABNT Catálogo}: confirmação da ABNT NBR 5419-1:2026 e consulta das edições indicadas.\\
\href{https://www.abntcatalogo.com.br/}{\texttt{abntcatalogo.com.br}}
\item \textbf{Equatorial Energia Maranhão}: lista vigente das normas de conexão e respectivos anexos, incluindo NT.00001 Rev.09, NT.00002 Rev.10, NT.00004 Rev.08, NT.00020 Rev.06, NT.00030 Rev.03 e NT.00042 Rev.04.\\
\href{https://ma.equatorialenergia.com.br/institucional/leis-e-normas-tecnicas/normas-tecnicas/instalacao-de-padrao/}{\texttt{Equatorial MA — Normas Técnicas de Conexão}}
\item \textbf{Ministério do Trabalho e Emprego}: NR-10 vigente até 31/05/2027 e nova redação com vigência em 01/06/2027.\\
\href{https://www.gov.br/trabalho-e-emprego/pt-br/acesso-a-informacao/participacao-social/conselhos-e-orgaos-colegiados/comissao-tripartite-partitaria-permanente/normas-regulamentadora/normas-regulamentadoras-vigentes/norma-regulamentadora-no-10-nr-10}{\texttt{MTE — página oficial da NR-10}}
\end{itemize}
\begin{caixaAt}
Consulta realizada em 13/08/2026. A lista da distribuidora pode mudar; no início de cada projeto, conferir novamente revisão, anexos e período de transição. Este material resume critérios e não reproduz nem substitui o texto integral das normas ABNT.
\end{caixaAt}
\end{frame}

\begin{frame}{Referências para a loja de shopping}
\small
\begin{itemize}
\item \textbf{ABNT NBR ISO/CIE 8995-1}: iluminação de ambientes de trabalho interiores; confirmar edição vigente no \href{https://www.abntcatalogo.com.br/}{Catálogo ABNT}.
\item \textbf{ABNT NBR 10898}: sistema de iluminação de emergência; usar em conjunto com o projeto de segurança contra incêndio aprovado.
\item \textbf{CBMMA}: \href{https://cbm.ssp.ma.gov.br/unidades-bm/unidades-administrativas/dat/legislacao-tecnica-dat/}{legislação técnica oficial}, incluindo NT 18 para iluminação de emergência e demais normas aplicáveis ao shopping.
\item \textbf{Manual técnico do empreendimento}: ponto de entrega interno, demanda concedida, proteção, medição, materiais, procedimentos de obra e comissionamento.
\item \textbf{Dados fotométricos e fichas dos fabricantes}: arquivos IES/LDT, fluxo, distribuição, potência, IRC, CCT, compatibilidade de controles, torque e manutenção.
\end{itemize}
\begin{caixaAt}
As metas numéricas do estudo da loja são premissas didáticas declaradas. O executivo deve confrontá-las com as edições normativas vigentes, a atividade visual real, a marca e as exigências do shopping.
\end{caixaAt}
\end{frame}

\section{Fechamento}
\begin{frame}{Síntese}
\small
\begin{caixaProjeto}[title=Competência de projeto]
O dimensionamento elétrico não é uma sequência isolada de fórmulas. Um projeto consistente coordena \textbf{carga, demanda, condutor, proteção, queda de tensão, aterramento, segurança, documentação e operação} como partes de um único sistema.
\end{caixaProjeto}
\vspace{2mm}
\begin{caixaObs}
Os exemplos são didáticos. Projetos executivos devem utilizar dados reais da instalação, critérios normativos vigentes, informações da concessionária e responsabilidade técnica profissional.
\end{caixaObs}
\end{frame}



\end{document}
